% NO \documentclass, NO \begin{document}, NO \end{document} here!
\chapter{Mediating the State: Middle Bureaucracy and Social Change}
\setcounter{section}{0}
\setcounter{subsection}{1}
\renewcommand{\thesection}{\arabic{section}}

\bigskip
\noindent\textbf{Part A}\medskip

\noindent\textbf{When Are Weak States \emph{Not} Weak?}
\noindent Across the global South, bureaucracies routinely dismissed as patronage-ridden or inert sometimes deliver services that rival far richer polities. Fortaleza treats 94 percent of households for sewage while Belém—operating under the same Brazilian federal charter—treats only 13 (World Bank 2018). Surabaya’s bus-rapid-transit network in Indonesia runs at 95 percent punctuality; Medan’s sister system limps along at 60 (ITDP 2019). South Africa’s Cape Town supplies water round-the-clock while Johannesburg imposed rationing in the very same drought year (McDonald & Ruiters 2005). Such contrasts are not statistical curiosities; they are patterned anomalies that expose blind spots in prevailing theories of state capacity. This question of variation in public goods in similar jurisdictions has long puzzled scholars of comparative politics. Around the world, one can find adjacent countries, regions, or even cities that share common histories or formal structures yet deliver strikingly different levels of public goods and services. A classic example comes from Italy: after new regional governments were established uniformly across the country, the northern regions flourished while the south lagged, a divergence that Robert Putnam 
\citeyear{putnam_making_1994} famously attributed to deeper variations in civic engagement and social capital.

Such contrasts are not confined to any one place. From Latin America to South Asia, we see cases where two jurisdictions starting with comparable institutional blueprints end up on very different trajectories of governance—one managing to provide reliable public services and maintain accountability, the other mired in dysfunction. These stark inconsistencies defy the expectation that similar starting conditions should produce similar outcomes. They beg a deeper exploration: when formal institutions are held constant, what explains why one locality succeeds in delivering basic services while another fails?

Popular explanations offer only partial answers. Modernization theory points to socioeconomic endowments, yet poorer Surabaya outperforms wealthier Medan. Developmental-state arguments emphasize cohesive Weberian bureaucracies, but India’s all-India services are uniform across Uttar Pradesh. New-institutionalist accounts highlight formal rules, yet Belém and Fortaleza share the same federal charter and procurement code. The disjuncture between theory and fact demands that we rethink the micro-mechanisms through which state capacity is constructed—and where it collapses.

This broad puzzle finds a vivid echo in urban India. In fact, India’s cities present a striking controlled comparison in governance. Unlike many federal countries where cities operate under different legal regimes, Indian cities share an almost identical institutional framework. The Indian Constitution (through the 74th Amendment) grants all cities the same formal status, and civil service rules and fiscal arrangements for municipal governments are largely standardized nationwide. Yet despite this common institutional template, basic service delivery varies enormously from city to city. Consider two large state capitals, Bhopal and Lucknow. Both are important historical cities and administrative centers in their respective states, and both operate under the same national urban governance laws and civil service structures. Yet Bhopal in recent years has been celebrated as one of India’s cleanest, best-serviced cities – in 2022 it climbed to become the 6th cleanest city in the country – while Lucknow until recently struggled with chronic service failures, ranking as low as 269th in cleanliness in 2017.  Although Lucknow has since improved its standings through focused efforts, the fact remains that for years these two cities – alike in formal powers and constraints – experienced vastly different governance outcomes. Such divergences are not limited to cleanliness rankings. 

Such divergences are not unique to India's cleanliness metrics but appear across diverse contexts and performance indicators worldwide. Consider two Colombian cities, Medellín and Cali, which illustrate how identical institutional frameworks can yield radically different social outcomes. Both are large urban centers operating under Colombia's decentralized governance system with identical legal authorities. Yet Medellín has transformed from the world's most dangerous city in the 1990s to a model of social innovation, implementing successful education programs, public libraries, and cable cars connecting poor neighborhoods to city services. Meanwhile, Cali despite having similar resources and formal powers continues to struggle with high violence rates, failing education systems, and limited access to public services in marginalized communities. Although Cali has recently attempted reforms, for decades these two cities—sharing the same institutional blueprint—demonstrated profoundly different capacities to deliver social services and ensure citizen security. All this is occurring under a common institutional umbrella of laws, bureaucratic rules, and fiscal formulas that apply to every city. 

The puzzle, then, is clear: how can regions, cities, districts and states that look institutionally identical on paper perform so unevenly in practice? What accounts for these “pockets of effectiveness” in a landscape otherwise marked by weak state capacity?

\textbf{Limitations of Existing Explanations}

To approach this puzzle, it is useful to consider three influential strands of theory that have tried to explain why some polities govern better than others: modernization theory, developmental state theory, and institutionalist theory. Each offers important insights; yet, as will be explained, each also has blind spots that limit its ability to fully account for the Indian urban divergence.

Modernization theory provides a classic starting point. Modernization scholars argue that as societies become more economically developed and socially complex, their political and administrative institutions also modernize, leading to better governance and public goods provision. In Seymour Martin Lipset’s famous formulation, “the more well-to-do a nation, the greater the chances it will sustain democracy”– and by extension, prosperity breeds the conditions for effective, accountable government (Lipset 1959). Early modernization theorists saw economic growth, education, and urbanization as driving forces that transform social structures and ultimately yield more capable governance and more responsive public institutions. However, the modernization perspective, with its emphasis on broad societal progress, often struggles with scale: it excels at comparing nations at different stages of development, but falters in explaining why units at a similar level of development diverge. In the context of Nigeria, modernization theory would predict roughly similar outcomes; yet, within one country, colonial-era divergences produced starkly different modernization outcomes. Under British rule, Western education and Christianity spread rapidly in the South but were strongly resisted in the Muslim North, where indirect rule left local structures intact. As a result, by independence, the South had far higher schooling and social indicators, while development in the North was “much slower,” creating a large internal gap. Modernization theory’s expectation of a uniform national trajectory falters here – cultural and institutional actors (missionaries, local rulers) created regional paths. The theory’s blind spot is failing to account for sub-national context: two regions under the same nation-state experienced the 20th century very differently, a pattern modernization theory alone wouldn’t predict.

Developmental state theory shifts the focus from social structure to the state apparatus. This school of thought, inspired by the spectacular success of East Asian countries in the late 20th century, argues that developmental outcomes depend on having a cohesive, capable bureaucracy and strategic state intervention in the economy. Scholars have credited a model of state-led capitalism called the “developmental state” with producing the economic miracles of Japan, South Korea, and Taiwan(Lee, 2020). In these cases, centralized governments with meritocratic bureaucracies and a strong sense of national mission engineered rapid industrialization and public service expansion. Key features of the developmental state are often said to include a stable, centralized authority, a highly trained bureaucracy with considerable autonomy, and an overarching commitment to developmental goals. At first glance, one might assume that the flip side of this argument is that countries without such state structures – often in South Asia or Africa – are doomed to poor governance. India, famously, has never been categorized as a classic developmental state; its postcolonial trajectory has been one of a “weak-strong state” – a democracy with pockets of capacity but also fragmentation and clientelism. Yet even within generally non-developmental states, we sometimes observe islands of good governance. The developmental state paradigm has an actor blind spot when it comes to these sub-national variations. It tends to treat the state as a unitary whole (either you have a cohesive high-capacity state or you don’t) and focuses on top-down institutions and national political leadership. What it overlooks are the actors operating below the apex of state power – the mid-level bureaucrats, municipal technocrats, and local political entrepreneurs who might drive change in localized settings even when the national state apparatus is imperfect. Developmental state theory, with its macro-institutional lens, has difficulty explaining why – Ceará, a relatively poor state in Northeast Brazil, became a role model for development through innovative governance(Tendler, 1997). “Despite having the 5th lowest GDP per capita” among Brazil’s 26 states, Ceará achieved “the largest increase in the national education quality index” since 2005. It did so by pioneering results-based financing in education, pushing universal literacy at the municipal level, and efficiently using scarce resources. Meanwhile, wealthier states like Rio or São Paulo did not see such rapid gains. This subnational success exposes a scale blind spot in developmental state theory: a state/provincial government can act as a developmental driver even if the national state isn’t uniformly effective. Ceará’s case also highlights actors beyond the central government – local officials and communities – as key to development, an aspect classic developmental state theory may undervalue.

Institutionalist theory offers another influential perspective, one that emphasizes the rules of the game. Both the new institutional economics and historical institutionalism have argued that the quality of governance and development outcomes is fundamentally shaped by the strength and design of institutions – the constitutions, laws, and formal organizations that structure incentives. In the grand comparisons of why nations succeed or fail, institutionalists like Douglass North, or Acemoglu and Robinson, assert that inclusive, well-designed institutions create prosperity, whereas extractive or dysfunctional institutions perpetuate poverty and misrule.A nation with secure property rights, rule of law, and accountable governance structures is expected to outperform one with corrupt or ineffective institutions. At first glance, the Indian city context seems like a perfect test of institutionalist predictions: here we have a set of cities with virtually identical formal institutions – the same legal mandates for city councils, the same civil service recruitment through national and state exams, the same fiscal transfer formulas and auditing mechanisms. If formal institutional design alone were the key determinant of performance, we would expect Bhopal, Lucknow, and other cities to deliver public services at roughly similar levels of quality. That they do not suggests that institutional theory, at least in its static form, is missing something. One issue is what we might call a temporal blind spot: institutionalist accounts often take a snapshot view of formal rules, failing to account for how the history and timing of reforms, or the informal institutions and norms overlaying formal structures, affect outcomes. In India, for instance, all cities received a similar institutional framework in the early 1990s reforms – but some cities might have had prior traditions or leadership cultures that allowed them to take advantage of those reforms, while others did not. Moreover, institutionalist theories can be actor-blind in their own way, treating institutions as if they operate automatically. In reality, institutions are animated by people. The same municipal act and service rules can be implemented in earnest in one city or ignored in another, depending on the initiative of local officials or the pressures of local politics. In fact, mounting evidence from governance research indicates that what truly drives public-sector performance are the political and organizational dynamics behind the scenes – the degree of elite commitment, the informal enforcement of rules, the networks of accountability – “rather than … formal political institutions themselves”(Hickey, 2023) Thus, a purely institutionalist theory cannot, by itself, elucidate why two cities with the same formal constraints end up miles apart in outcomes. We must look beyond the letters of the law to the living process of governance.

In summary, each theoretical tradition contributes a piece to the puzzle but leaves important gaps. Modernization theory reminds us that social context and development levels matter, but it glosses over local deviation and agency. Developmental state theory underscores the importance of a capable bureaucracy and state commitment, but it usually operates at the wrong level of analysis for intra-country variation and neglects the role of individual actors apart from top leaders. Institutionalist theory properly highlights the significance of rules and structures, yet by assuming identical rules should yield identical results, it fails to explain the nuanced divergence among Indian cities. Each suffers from one of the blind spots identified above – be it scale, actors, or time. To solve the puzzle of variation in public services in weak states, we need an explanation that operates at the right scale (the city/meso level), that recognizes the agency of key actors, and that accounts for the temporal development of governance practices. The next section introduces the core argument of this dissertation, which aims to fill these gaps.

\textbf{Conditional Autonomy & Meso Level Bureaucrat: A New Explanation}
This dissertation argues that the missing piece in understanding why some jurisdictions placed under the same scope condition of weak outperform others lies in a concept I call \textbf{conditional autonomy} – essentially, the presence of \textit{negotiated professional discretion} at the \textbf{meso-bureaucratic level}. The central thesis is that in low-capacity, patronage-driven environments, governance outcomes hinge on the dynamic between political leaders and bureaucratic officials. Granting mid-level bureaucrats conditional autonomy—some independent authority even amid patronage pressures—leads to significantly better public service delivery than a fully collusive politician–bureaucrat nexus. This performance gap is most pronounced in complex, networked public goods (such as urban water systems or electricity grids) that require coordinated administration and cannot be parceled out as simple patronage rewards. In contrast, purely collusive systems fragment services into localized patronage fiefdoms, and the ability (or failure) to deliver broad, non-targetable services starkly reveals the underlying strength of bureaucratic autonomy.

In plainer terms, jurisdictions achieve better governance outcomes when meso-level public officials (neither the high-ranking elite nor the street-level workers) are able to carve out a sphere of autonomy to act professionally and effectively, and when that autonomy is tolerated – even encouraged – by political leaders \textit{under certain conditions}. It is in these spaces of conditional autonomy that “pockets of effectiveness” emerge. We can define a pocket of effectiveness as a public organization or jurisdiction that succeeds in delivering public goods above the prevailing standard of its environment \citep{hicken_building_2009}. Scholars have noted such pockets in otherwise weak governance settings – for example, a well-functioning electricity board or health department that stands out in a country known for bureaucracy and corruption. These pockets of effective performance, by their very existence, “contradict the sweeping generalizations” (Hickey, 2023)about overall state failure and demand an explanation at a finer grain.

The concept of \textbf{conditional autonomy} provides that explanation by focusing on the relationship between bureaucratic agents and their political principals. Autonomy here does not mean complete independence or a legal grant of freedom; rather, it is \textit{conditional and negotiated}. Talented bureaucrats gain a degree of professional discretion – the freedom to make decisions, innovate, and enforce rules in the public interest – when political leaders see it as beneficial (or at least not harmful) to their own goals. This often involves a tacit bargain. Bureaucrats are given leeway to improve services and implement reforms, and in return, they deliver results that elected officials can claim credit for, or at least they avoid embarrassing failures that would anger the public. Political leaders, on their side, agree (implicitly) to not micromanage every decision or subvert the bureaucracy for patronage, so long as the bureaucrats’ actions align with certain political considerations (for instance, not undermining key supporters or jeopardizing re-election prospects). The discretion is conditional – it can be rescinded if the bureaucrat oversteps or if political calculations change – but it can be enough to permit meaningful improvements in governance. Crucially, this arrangement hinges on the professional norms and initiative of the bureaucrats (their willingness and ability to use the freedom productively) and on the strategic calculations of politicians (their willingness to prioritize broad service gains, at least temporarily, over narrow patronage or control).

By focusing on these meso-level interactions, the dissertation locates agency and reform capacity in a place that broader theories rarely probe. Conditional autonomy suggests that even in a formally uniform and overall weak institutional context, \textit{variation can arise from how local actors navigate and negotiate the implementation of those institutions}. One city becomes an over-performer because at some point a critical mass of bureaucratic officials and political actors arrived at an understanding that enabled effective problem-solving. Another city remains an under-performer because such an understanding never took hold – perhaps local officials were more tightly shackled by political interference, or less inclined to push boundaries, or political incentives favored maintaining the status quo. 

This theory of conditional autonomy does not discard structural factors – indeed, it operates within the constraints of formal institutions and socio-economic conditions – but it zeroes in on how bureaucratic agency and political strategy interact at the city level to create pockets of good governance inside a generally challenged system. It acknowledges, as recent research has emphasized, that high performance in the public sector often requires political buy-in to shield agencies from corrosive pressures \citep{klitgaard_high-performance_2005}. At the same time, it insists that bureaucrats are not just passive instruments of development or decline; they are players who can sometimes negotiate the terms of their engagement. By explaining pockets of effectiveness as a product of negotiated space for bureaucratic professionalism, the argument bridges the gap between the macro and micro, showing how local agency can modulate the outcomes predicted by broader structural theories.

To investigate this argument, the dissertation adopts a comparative case study approach, focusing on two emblematic cities: Bhopal and Lucknow. These cases are selected as a “most similar” pair – similar enough in formal institutional arrangements and other background factors to make their divergent outcomes puzzling, yet different in performance in a way that begs for explanation. Both Bhopal and Lucknow are capitals of major Indian states (Madhya Pradesh and Uttar Pradesh, respectively). Both host a full array of government agencies and enjoy the powers and responsibilities conferred to all municipal corporations in India. Both grapple with the typical challenges of Indian urban centers – rapid growth, infrastructure deficits, fiscal constraints – and both operate under the same national legal framework for urban governance. By holding constant many high-level factors (constitutional status, administrative rules, etc.), these cases allow us to zoom in on the meso-level dynamics that differ between them. In essence, Bhopal and Lucknow provide a test of what happens when the software of governance (the informal arrangements and leadership practices) differs, even though the hardware (the formal institutions) is the same.

The research design weaves together multiple methods to build a robust account of governance trajectories in each city. At its core lies a controlled within-city comparative design that systematically pairs neighborhoods across Bhopal and Lucknow. Four wards in each city were selected using matched sampling methodology, creating comparable neighborhood pairs—a low-income slum-dominated ward in Bhopal matched with its demographic equivalent in Lucknow, for instance. This micro-level matching allows "like with like" comparison, isolating governance effects from demographic variation.

Archival research provides the historical foundation, drawing on decades of government orders, city council proceedings, committee reports, and local press archives. These sources reveal critical junctures—when sanitation initiatives launched or administrative powers shifted—capturing moments when conditional autonomy emerged or collapsed. Complementing these official records, extensive elite interviews with current and former administrators, elected officials, and knowledgeable observers illuminate the informal negotiations and power dynamics that rarely leave a paper trail. These conversations often revealed how municipal commissioners negotiated with chief ministers for project autonomy, or how political pressures operated behind closed doors.

The research also employed deep ethnography among municipal employees, embedding with frontline workers to document their daily interactions and decision-making processes. This ground-level immersion revealed how formal rules translate into practice, exposing the micro-mechanisms through which bureaucratic autonomy operates—or fails to operate—in each city. A face-to-face comprehensive household survey of 3,754 residents, equally divided between cities, measured citizens' experiences with public services and their patterns of seeking redress, confirming measurable differences in service delivery while revealing variations in bureaucratic responsiveness.

Finally, administrative datasets comparing municipal budgets, staffing records, and service metrics help identify whether differential spending, reform timings, or institutional changes explain the cities' divergent outcomes. Through this triangulation of archival, interview, ethnographic, survey, and administrative data, the research reconstructs each city's governance trajectory—pinpointing when changes occurred, who drove them, and what effects they produced.

It is important to note that while the core of the research is this in-depth comparison of Bhopal and Lucknow, the dissertation also reflects on broader evidence from other Indian cities where relevant. The aim is not merely to tell two idiosyncratic stories, but to use these cases to shed light on a general phenomenon. The processes identified – such as how a reformist bureaucrat can form a coalition with civil society to pressure politicians, or conversely how political patronage networks can strangle administrative initiative – have analogues in other cities. Thus, the findings from Bhopal and Lucknow will be used to suggest a framework applicable to understanding success and failure in urban governance across India and, potentially, in other national contexts where formal institutions are shared but outcomes diverge.

By answering this puzzle of jurisdictions with identical formal institutions produce radically different governance outcomes, this dissertation makes several wider contributions to political science and public policy scholarship.

First, the project refines bureaucratic theories by introducing the concept of conditional autonomy—a negotiated equilibrium that emerges in contexts where neither idealized Weberian administration nor pure patronage prevails. Traditional bureaucratic theory frames autonomy in stark terms: shielded rule-bound apparatus or captured instrumentality. Yet evidence from Bhopal demonstrates how bureaucrats negotiate a dynamic "gray zone," securing sufficient insulation for technical standards while accommodating political imperatives. This conceptual innovation moves beyond binary framings to capture how meaningful reforms unfold within politically constrained environments. The significance lies in demonstrating that governance improvements need not await wholesale institutional transformation; rather, partial autonomy can yield incremental yet substantive gains that accumulate over time.

Second, the project advances state capacity theory by elevating mid-level bureaucrats from analytical obscurity to theoretical centrality. While existing scholarship fixates either on elite policymakers with brief tenures or street-level workers inundated with operational demands, this work illuminates the "missing middle"—state-level officials who function as critical institutional intermediaries. Building on embedded autonomy frameworks \citep{evans_bureaucracy_1999}\citep{migdal_strong_1988}\citep{tsai_accountability_2012}, I demonstrate how these actors translate policy abstractions into operational realities, orchestrate horizontal coordination across fragmented agencies, and forge vertical linkages between strategic objectives and implementation. Their role becomes particularly salient in India's paradoxical state structure—simultaneously "weak-strong" \textbf{(Rudolph & Rudolph 1987)}, "failed developmental" \textbf{(Herring 1999)}, and "flailing" \textbf{(Pritchett 2009)}—where mid-level bureaucrats transform sporadic reforms into sustained capacity-building.

Third, the dissertation enriches urban governance scholarship by demonstrating how incremental progress depends critically on bureaucratic intermediation. While recent work recognizes that gradual institutionalization enables enduring improvements \citep{eduardo_politics_2021}, my comparative findings reveal a more nuanced mechanism: conditional autonomy among embedded bureaucrats allows them to navigate patronage-laden environments and translate formal reforms into tangible outcomes. This insight challenges prevailing models that attribute incremental gains solely to stable political coalitions or municipal oversight structures. Instead, I show how bureaucratic agency shapes contingent achievements—where officials strategically adapt, selectively insulate, and incrementally steer systems toward more reliable public goods provision.

Fourth, the project makes a significant methodological contribution by transcending common limitations in governance research. Rather than relying on static cross-sectional analyses or narrow principal-agent frameworks, this study employs longitudinal process tracing to capture bureaucratic adaptation over time \cite{cox_electoral_1986}

(Weingast, 1984; McCubbins et al., 1987). By following how officials respond to shifting political priorities, evolving donor agendas, and infrastructural investments across extended periods, the research reveals both trigger mechanisms and accumulation processes. This temporal approach explains not only what initiates episodes of fragmentation or cooperation, but how incremental adjustments crystallize into durable patterns of institutional capacity—or persistent deadlock.

These contributions collectively reframe our understanding of state capacity formation. Rather than conceptualizing capacity as uniformly distributed or achievable only through top-down reform, this work demonstrates how pockets of effectiveness emerge through sustained negotiation between political constraint and bureaucratic agency. By centering analysis on conditional autonomy and mid-level intermediation, the dissertation provides both theoretical refinement and practical insights for nurturing state capacity from within existing constraints.

\section{Three Blind spots in State Capacity Scholarship}

As outlined in the previous section, existing theoretical frameworks fail to explain subnational
governance disparities. 

Modernization theory cannot account for wealthy regions that underperform
poorer ones, institutionalism wrongly assumes identical formal rules yield identical outcomes, and
developmental state theory operates only at the national level, missing internal variation. Empirical
evidence confirms these limitations: standard national-level variables leave subnational differences
largely unexplained. This theoretical gap demands new approaches that conceptualize state capacity as
a variable, emergent property within countries rather than assuming uniform national outcomes.

Why have our dominant theories missed the mark on subnational governance variation? This dissertation
identifies three key blindspots in prevailing theories of state capacity and development:

\textbf{Scale Error}: The first blindspot is a mismatch of scale. Modernization and institutionalist theories tend
to treat the nation-state as the unit of analysis, implicitly assuming a homogenous polity. This nationallevel
focus overlooks crucial within-country differences – what O’Donnell called the “unevenness of
stateness” across territory. By aggregating to a single national value for “state capacity” or
“governance,” scholars commit a scale error, conflating the national formal structure with the on-theground
reality in every locality. The result is that country-level scores hide enormous within-country
variability. For instance, Freedom House might rate a country as “Partly Free,” but within that country
some provinces are virtually “Not Free” in practice while others approach the ideal of liberal
governance. Any theory that treats state capacity as a binary national attribute (“weak” or “strong” state)
fails to capture this spatial variation. The scale error thus leads to misleading conclusions – we might
credit institutional successes to a country’s laws or resources, when in fact only select regions are driving
the outcomes (or vice versa, national failures might obscure local successes).

\textbf{Actor Blindspot}: The second blindspot is an oversimplification of who actually drives governance
outcomes. Many theories are actor-blind beyond abstract categories like “elites” or “the state.”
Modernization theory focuses on broad social classes and income levels; institutionalism focuses on
rules of the game; developmental state theory focuses on top leaders and high-level bureaucratic bodies.
Missing from these accounts are the meso-level actors – the middle-tier officials, local administrators,and front-line bureaucrats who actually implement policies day-to-day. The importance of these actors
has been well documented in public administration research: Michael Lipsky observed that “public
service workers, in effect, function as policy decision makers, as they wield their considerable discretion
in the day-to-day implementation of public programs.” The cumulative effect of street-level and midlevel
decisions can “reroute the intended direction of policy,” often producing outcomes very different
from what central laws envisioned. In short, policy on paper is not the same as policy in practice, because
human agents intervene at every step. The mainstream comparative theories exhibit an actor blindspot
by assuming that either rules or resources translate automatically into outcomes. They rarely account
for how bureaucratic agency, motivation, and behavior mediate (or frustrate) the link between
institutions and results. Ignoring these implementers is perilous: even a well-designed institution can
falter if those charged with carrying it out lack incentives or capability; conversely, a flawed institutional
setup can be salvaged by creative, committed bureaucrats. The actor blindspot thus contributes to the
failure of dominant theories to predict subnational divergence – because where governance is better,
often the difference lies in who is doing what, not just in what institution they operate.

\textbf{Temporal Mis-specification}: The third blindspot involves the temporal dimension of state capacity.
Classic theories often mis-specify how capacity develops or changes over time. Modernization theory
envisioned a more or less automatic, linear progression – as economies grow and societies become
complex, political development (including state capacity) is supposed to follow. Developmental state
theory emphasizes critical junctures and big pushes (e.g. Meiji Restoration in Japan, the Park Chunghee
era in Korea) where capacity is rapidly built by a visionary leadership; institutionalist analyses often
focus on founding moments (constitution-writing, regime change) or path-dependent persistence. What
these approaches share is a relative neglect of incremental and iterative change. They have trouble
explaining how two regions that started with identical institutions gradually diverge. If one assumes
capacity is set by historical structural factors (e.g. colonial institutions or initial resource endowments),
then subnational divergence today would trace to some distant divergence in the past – an implausible
story in many cases. Instead, evidence suggests that institutions “change in subtle and gradual ways over
time,” and such slow, piecemeal changes can be equally consequential for political outcomes as dramatic
ruptures. The temporal mis-specification of conventional theories is that they either expect immediate
effects from formal institutions (ignoring the time it takes for new rules to embed, or for old practices
to erode) or they expect change only from large exogenous shocks. In reality, governance capacity often
emerges cumulatively: small policy experiments, marginal budget increases, one competent appointment
at a time. By treating state capacity as static or as a binary that flips when a threshold is crossed,
prevailing theories miss the ongoing, negotiated process of change. This dissertation argues that we must
pay attention to process – the temporal sequence by which certain subnational units inch ahead of others,
even under the same overall institutional environment.
In summary, the failure of dominant approaches to explain uneven subnational governance is rooted in
these blindspots. They looked too high up (the national scale), too far away (abstract structures over
human agents), and too briefly (snapshots or dramatic moments over gradual evolution). How, then, can
we do better? Addressing the puzzle requires a theoretical framework that corrects these blindspots –
one that zooms into the subnational scale, puts key actors under the microscope, and respects the
temporal dynamics of change.

\textbf{Meso-Level Bureaucrats and Conditional Autonomy: A New
Perspective}

the often-neglected meso-level of governance: the provincial department chiefs, district collectors, city
managers, program directors – and their field staff – who operate between high-level policy and street
level implementation. These meso-level bureaucrats are the connective tissue of the state. They translate
broad policies into local practice, allocate resources on the ground, and interact with both political
superiors and community actors. C\textit{rucially, they also possess a degree of autonomy in their domain. The
argument advanced here is that variation in subnational governance capacity largely stems from how
meso-level bureaucrats exercise what I call “conditional autonomy.”}


Conditional autonomy refers to a bounded, context-dependent form of bureaucratic independence.
Unlike the robust Weberian autonomy celebrated in classical bureaucratic theory – where officials are
insulated by meritocratic norms and can execute policy without interference – conditional autonomy is
fragile and contingent. Bureaucrats have some leeway to innovate or enforce rules diligently, but that
leeway can quickly evaporate if political circumstances shift or if they overstep informal limits.
Independence is thus “conditional” on factors such as support from certain politicians, a tacit bargain
with local elites, or staying within the radar of central authorities. Recognizing this nuance is vital. In
many developing contexts, we will not find perfectly autonomous bureaucracies (indeed, attempts to
create them often fall prey to patronage or politicization). What we find instead are pockets of autonomy
– specific departments, agencies, or local offices where officials manage to carve out space to get things
done, even if that space is precarious.

Why focus on these meso-level bureaucrats and their conditional autonomy? Because they offer a
mechanism for incremental governance improvements that does not require an overhaul of formal
institutions. When mid-level officials take initiative – say, a district health officer campaigns for
immunization with unusual zeal, or a municipal administrator streamlines a permit process – they can
improve outcomes at the margin. Over time, a series of such incremental changes can aggregate into a
substantial divergence in capacity between two otherwise similar regions. One region’s officials, given
slightly more autonomy to solve problems (perhaps because a reformist minister empowered them, or
local civil society held them accountable, etc.), will gradually build better performance metrics – higher
tax collection, lower infant mortality, more reliable utility services. The neighboring region, with more
constrained or demotivated officials, stagnates. Small differences in bureaucratic practice compound
into large differences in governance. This bottom-up, agent-driven path to capacity stands in contrast to
the top-down, structural path assumed by dominant theories.

Concrete examples illustrate the point. As discussed briefly in the introduction, in Brazil’s Ceará state
– one of the country’s poorest – a series of innovative programs in the late 1980s resulted in dramatic
improvements in public health and administration, even though Brazil’s national institutions were
unchanged. Studies of this “Ceará miracle” found that state-level technocrats empowered local health
workers through new accountability and incentive structures, essentially granting them conditional
autonomy to respond to community needs. Mayors could no longer meddle in hiring these workers, and
the state government closely monitored local delivery while enlisting civic groups, creating a balance
of oversight and freedom. The result was a surge in vaccinations and a drop in child mortality – outcomes
far better than in other equally poor Brazilian states. The Ceará case underscores how mid-level
bureaucratic initiative, under the right conditions, can defy the expectations of institutional determinism.
Similarly, in India, researchers have noted that some district collectors (the chief administrative officers
in a district) manage to significantly improve local project implementation by creatively using their
limited discretionary funds or by coordinating between siloed departments, whereas others simply
follow routine. The difference often comes down to whether the bureaucrat feels politically supported
and professionally motivated to take risks – a conditional autonomy conferred by an enabling
environment.

What these examples share is a vision of change driven from the meso level. Rather than waiting for a
sweeping national civil service reform or a new law from parliament, improvements emerged from
bureaucrats “in the middle” making the most of what they had. This perspective directly addresses the
earlier blindspots: it operates at the appropriate scale (the subnational unit, where variation occurs), it
highlights specific actors (bureaucrats as change agents or obstacles), and it embraces the temporal
accumulation of reforms (capacity as emergent through repeated actions). By focusing on meso-level
bureaucrats, we also acknowledge that governance is not monolithic – \textit{within the same institutional
framework, some parts of the bureaucracy can perform well while others fail}. The task is to explain why
and how those differences arise.

One key insight is that meso-level bureaucrats act as intermediaries who must negotiate on two fronts:
vertically with higher authorities (politicians, senior officials) and horizontally with local actors
(citizens, businesses, interest groups). Their autonomy is won and used in these negotiations. A district
education officer, for instance, might negotiate with the provincial ministry for flexibility in allocating
teachers, while simultaneously negotiating with village leaders to cooperate on monitoring school
attendance. If either negotiation breaks down – if the ministry clamps down with strict directives, or if
the community mistrusts the officer – her room for maneuver shrinks. But if she secures a tacit approval
from above and buy-in from below, she can push reforms that stick. In essence, state capacity at the
local level is co-produced by these ongoing negotiations between state and society, and between
different levels of the state hierarchy.

\textbf{State Capacity as Emergent and Negotiated (Not Given or Taken)
}Stepping back, this approach suggests a broader reconceptualization: state capacity is not a trait that a
state simply “has” or “lacks” at the national level – it is an emergent, negotiated phenomenon that often
crystallizes unevenly across subnational spaces. It emerges from countless interactions: directives from
the center, responses at the periphery, innovations tried, responses given, bargains struck (or missed)
over resources and authority. In this view, a “strong state” is not one that uniformly penetrates its
territory with authority (as classical theories might assume); rather, a state’s strength varies within, and
pockets of effectiveness must coalesce through practice and contestation.

Consider again the notion of emergence: like a complex system, governance capacity results from the
aggregation of local actions and adaptations. No central planner can design the perfect local governance
for every village – instead, as local officials and communities solve problems in their own contexts,
some patterns of effective governance emerge, which can then spread or be scaled up. This bottom-up
emergence means that the path to high state capacity is rarely uniform or centrally orchestrated. National
indicators may improve only when enough local units have transformed their practices – which might
happen in one region decades before another. Consequently, any theory that treats state capacity as
homogeneous is bound to miss where the real changes happen.

Likewise, negotiation is at the heart of building capacity. This is not negotiation in the narrow sense of
formal bargaining, but in the broader sense of constant alignment-seeking among actors with partially
conflicting interests. Meso-level bureaucrats negotiate with political principals for support and
autonomy; central agencies negotiate with provincial authorities over policy implementation; local
officials negotiate with social groups over compliance and cooperation. Through these interactions, state
authority becomes “brokered” rather than simply imposed. Where governance improves, it is often
because a coalition or understanding was negotiated that allowed reforms to take root – for example, a
deal between a governor and bureaucrats that trading a bit of patronage for a lot of performance is
acceptable, or a social contract where citizens agree to pay taxes once they see services improve. These
negotiated arrangements can differ vastly from one locality to another, even under the same formal
institutional rules.

In contrast to a binary view (“State A is capable, State B is not”), the emergent-negotiated perspective
sees capacity as a spectrum – with each subnational unit somewhere on that spectrum, possibly moving
over time. We therefore get a more nuanced picture: a country might simultaneously exhibit 19thcentury
levels of bureaucratic dysfunction in one province and 21st-century best practices in another.
Instead of averaging them out, our theory seeks to explain that divergence. Importantly, this perspective
resonates with and builds on a growing body of research pointing to “pockets of effectiveness” in
otherwise weak states and the idea that even in broadly neo-patrimonial or low-capacity environments,
some institutions “work” due to specific political and bureaucratic arrangements. It also aligns with
historical observations that states rarely modernized evenly – for example, mid-20th-century scholars
of development noted that effective administration often began in select sectors (taxation, military)
before expanding outward. The contribution here is to formalize that understanding for contemporary
governance: state capacity must be understood as the cumulative result of localized, negotiated successes
(and failures), rather than as an exogenous national endowment.

To clarify how this dissertation’s approach diverges from the dominant theoretical traditions, Table 1
summarizes the key differences in focus and assumptions, as well as how the present theory addresses
their blindspots:

\begin{table}[H]
\centering
\caption{Contrasting Theoretical Approaches to State Capacity}
\small
\begin{tabularx}{\textwidth}{|p{3.5cm}|X|X|}
\hline
\textbf{Theoretical Approach} & \textbf{Core Assumptions about Governance} & \textbf{Blindspots and Limitations} \\
\hline

\textbf{Modernization Theory} \newline (Lipset; Rostow) & 
Socioeconomic development transforms society, leading to more effective governance and democracy as an almost natural progression. Governance capacity rises with wealth and education in a linear fashion. & 
Assumes a uniform national trajectory of progress. Fails to explain why regions at similar development levels govern differently. Treats state capacity as an automatic byproduct of growth, overlooking deliberate institution-building and local variation. \\
\hline

\textbf{Institutionalism} \newline (North; Acemoglu \& Robinson) & 
“Institutions rule”: Formal political and legal institutions (constitutions, electoral systems, bureaucratic structures) determine governance outcomes by shaping incentives and behavior. If rules are well-designed, good governance follows. & 
Assumes formal institutions have the same effect everywhere. Misses the role of implementation and informal practices. Actor blindspot: treats the state as a unitary actor, ignoring agency of mid-level officials. \\
\hline

\textbf{Developmental State Theory} \newline (Johnson; Evans; Kohli) & 
High state capacity is achieved by cohesive, centralized states that purposefully intervene in the economy. Strong bureaucratic institutions insulated from politics (often born out of critical junctures or elite vision) drive rapid development. & 
Focuses on the nation-state as the unit. Cannot explain why some subnational units thrive without a full “developmental state.” Overlooks incremental improvements. Assumes uniform capacity. \\
\hline

\textbf{Emergent-Negotiated Perspective} \newline (This Dissertation) & 
State capacity is an emergent property from cumulative local-level improvements. Meso-level bureaucrats with conditional autonomy incrementally solve problems. Effectiveness is co-produced through negotiations with political leaders and communities. & 
Embraces subnational variation as fundamental. Focuses on actors and process. Captures that capacity is partial, emergent, and spatially inconsistent. Explains why pockets of effectiveness appear under the same formal rules. \\
\hline

\end{tabularx}
\label{tab:theoretical-approaches}
\end{table}

\textbf{Conclusion and Chapter Outline}
In repositioning state capacity as something that happens rather than something a state possesses
outright, this chapter has set the stage for a new theoretical framework. We have argued that
traditional theories, while powerful in other contexts, overlook the subnational puzzle due to scale
errors, actor blindspots, and temporal mis-specifications. By contrast, a focus on meso-level
bureaucrats and their conditional autonomy allows us to see how governance can improve (or
stagnate) within the same institutional shell, through emergent and negotiated processes on the ground.

The next section will build directly on these ideas. It will formally theorize the role of meso-level
bureaucratic agency in state capacity, developing measurable concepts for “conditional autonomy” and
laying out hypotheses for how these mid-level actors drive incremental change. In doing so, sectopm
will also detail the research design for examining subnational variation in governance, including how
to identify and measure those elusive pockets of effectiveness. By moving from the broad narrative of
this introduction to a precise theoretical model, the coming section aims to operationalize the
dissertation’s core argument – that the key to understanding divergent governance outcomes lies in the
often invisible work of bureaucrats negotiating space to make the state work. Ultimately, this
theoretical grounding will guide the empirical analyses in subsequent chapters, which will test these
ideas across different regions and cases, showing when and where conditional autonomy leads to
tangible gains in governance capacity. Section 3 situates Indian cities as strategic sites for theorybuilding
in state capacity, examining what makes them theoretically significant cases. Section 4 then
delves into the empirical evidence through detailed comparative case studies of Bhopal and Lucknow.

In sum, the dissertation advances a comparative politics of the subnational state, one that
acknowledges that state capacity is not monolithic but built from the ground up. By recognizing the
centrality of meso-level bureaucrats in this process, we move closer to solving the puzzle of why

identical institutions can yield such divergent results – and how nurturing the right conditions for
bureaucratic autonomy and innovation can gradually turn weak governance into capable governance,
one district at a time.

\section{Theory and Measurement: Conditional Autonomy and
Meso Level Bureaucratic Action}

To bridge the theoretical gap between institutions and outcomes, this section develops a framework
centered on the relationship between political patronage and bureaucratic autonomy, demonstrating how
this dynamic determines governance performance through distinct modes of operation.

Much of the governance literature contrasts two ideal types: programmatic politics vs. patronage
politics, and autonomous bureaucracy vs. instrumental bureaucracy. In programmatic regimes,
politicians supposedly allocate resources based on universal programs, while in patronage settings, they
distribute benefits as particularistic favors. Similarly, the Weberian ideal envisions autonomous
bureaucracy governed by impersonal rules, while its opposite—often termed "neopatrimonial"—
represents complete bureaucratic instrumentalization. While the literature often terms this opposite as
"neopatrimonial," I use the term "collusion" to emphasize the active collaboration between political
patrons and instrumentalized bureaucrats in the particularistic distribution of resources. This
terminological choice highlights the relational nature of this governance mode—it requires both willing
political masters and compliant bureaucratic agents working together to distribute state resources for
private and political gain.

These dichotomies are analytically useful, but many real-world governance contexts do not fit neatly at
either extreme. Especially in developing and weak-capacity states, we observe a spectrum of
intermediate cases that combine elements of both logics. Figure 1 maps governance modes along two
dimensions: the horizontal axis measures patronage influence (from low on the left to high on the right),
while the vertical axis measures bureaucratic autonomy (from low at the bottom to high at the top). This
produces four distinct governance modes, each with its own internal logic and performance implications.

\begin{figure}[H]
    \centering
    \includegraphics[width=0.85\textwidth]{figures/01_1.png} % Replace with your filename
    \caption{Your descriptive caption goes here.}
    \label{fig:your-label}
\end{figure}

\begin{itemize}
    \item Top-left (Programmatic-Weberian): The idealized model of high bureaucratic autonomy and
low patronage influence, where professional civil servants implement universal policies.
 \itemTop-right (Conditional Autonomy): The critical intermediate zone of high bureaucratic
autonomy despite high patronage influence. Here, bureaucrats maintain professional discretion
while navigating political pressures. Bhopal exemplifies this mode.
 \itemBottom-left (Bureaucratic Authoritarianism): The relatively rare configuration of low
bureaucratic autonomy with low patronage influence, representing systems lacking both
direction and agency.
 \item Bottom-right (Collusion): The trap of low bureaucratic autonomy with high patronage
influence, where bureaucrats
\end{itemize}

The framework's key insight centers on the oval in the patronage half of the matrix—the empirical reality
for most developing countries. From this position, meso-level bureaucracy faces a crucial fork: either
ascend toward conditional autonomy or descend into collusion. This choice determines not just
governance quality but whether certain essential public goods can be delivered at all.

\begin{figure}[H]
    \centering
    \includegraphics[width=0.85\textwidth]{figures/01-01.png} % Replace with your filename
    \caption{Figure 1: Relationship between bureaucracy and Politics}
    \label{fig:your-label}
\end{figure}

Figure 2 presents the empirical distribution of countries using V-Dem data for 2023\footnote{The horizontal axis measures political patronage influence (using reversed 'v2psprlnks' where higher
values indicate more patronage), while the vertical axis measures bureaucratic autonomy (using
'v2clrspct' where higher values indicate greater autonomy). The quadrants are delineated by thresholds
at x=0 for patronage and y=0 for autonomy, corresponding to our theoretical typology}. The horizontal
axis measures political patronage influence, while the vertical axis measures bureaucratic autonomy.
The quadrants are delineated by thresholds at x=0 for patronage and y=0 for autonomy, corresponding
to our theoretical typology. In this graph:

\begin{itemize}
    \item Top-left quadrant: Countries like Germany, Switzerland, and the United States cluster here,
representing established programmatic-Weberian systems with low patronage and high
bureaucratic autonomy.
\item Top-left quadrant: Countries like Germany, Switzerland, and the United States cluster here,
representing established programmatic-Weberian systems with low patronage and high
bureaucratic autonomy.
\item Top-right quadrant: Countries like Singapore and post-conflict states occupy this space,
demonstrating that bureaucratic autonomy can coexist with patronage politics.
\item  Bottom-left quadrant: Relatively few countries occupy this space of low patronage and low
autonomy, suggesting this is an unstable or transitional configuration.
\item Bottom-right quadrant: Most developing countries cluster here, including India, Indonesia,
Ghana, and Afghanistan, confirming that collusive governance dominates in patronage-heavy
environments.

\end{itemize}

This clustering reveals that for most of the developing world, the critical governance challenge isn't
eliminating patronage—politically impossible in the short term—but preserving or building bureaucratic
autonomy within these constraints. The vertical spread on the right side demonstrates that bureaucratic
autonomy under patronage conditions is an important variable. This positioning fundamentally reframes
our understanding of governance possibilities. Rather than viewing patronage as uniformly detrimental,
the framework reveals how bureaucratic autonomy determines whether patronage contexts produce
functional governance or collusive dysfunction. The key distinction isn't the presence of patronage which is a given in most developing contexts—but whether bureaucracy maintains sufficient autonomy
to move above or below the Y axis in the two graphs that is does the bureaucracy get conditional
autonomy to innovate or it gets instrumentalized to maintain the status quo. That in essence is the point
of this dissertation.
Understanding these distinct governance modes—particularly the crucial distinction between
conditional autonomy and collusion—sets the stage for examining their concrete implications: how
these different bureaucrat-politician relationships affect the actual delivery of public goods and services.


\section{Public Goods Delivery: Why Conditional Autonomy Excels with Non-Targetable Goods}

The theoretical distinction between conditional autonomy and collusion becomes most consequential when examining public goods delivery, particularly through the critical divide between \textbf{targetable} and \textbf{non-targetable} goods that reveals why bureaucratic autonomy matters most for essential infrastructure and networked services (Table~\ref{tab:targetable-non-targetable}).

Targetable goods are those that can be organized within a district or jurisdiction without requiring network-wide coordination. Examples include ward-level road works, neighborhood solid waste collection, and installation of tubewells or handpumps. These services exhibit spatial divisibility and require action from only one or two state actors. Given their tightly bounded costs, timelines, and benefits, politicians can promise resurfaced lanes or improved waste collection within a single electoral cycle. Engineering staff can execute these works without extensive coordination. Crucially, politicians maintain complete discretion over where and when services are delivered, allowing strategic allocation for political gain.

In contrast, \textbf{non-targetable goods}—including piped water supply, sewerage-linked household toilets, and grid electricity upgrades—require network integrity, synchronized investment, and coordination among water boards, sewerage agencies, power utilities, and various levels of government. With gestation periods extending over years and benefits diffusing across boundaries, accountability becomes difficult to assign. Most importantly, the technical characteristics of these goods eliminate political discretion in distribution, making them less attractive for patronage-driven politics unless strong higher-tier leadership ensures coordination.

This asymmetry explains why jurisdictions frequently demonstrate rapid, visible improvements in road infrastructure and tubewells, while core network services like piped water remain underdeveloped. Politicians favor targetable goods because they offer maximum discretion, while non-targetable goods require durable coalitions and commitments that survive electoral transitions.

\begin{table}[H]
\centering
\caption{Targetable vs Non-Targetable Goods}
\label{tab:targetable-non-targetable}
\small
\begin{tabularx}{\textwidth}{|l|X|X|}
\hline
\textbf{Aspect} & \textbf{Targetable Goods} (roads, neighborhood trash, tubewells, handpumps) & \textbf{Non-Targetable Goods} (piped water, sewer-linked toilets, grid electricity) \\
\hline
Spatial divisibility & High – can be delivered ward-by-ward or even lane-by-lane & Low – networks must function city-wide to be viable \\
\hline
Administrative actors & One or two municipal departments; minimal external coordination & Multiple utilities and regulators; cross-departmental synchronization essential \\
\hline
Gestation period & Short (weeks–months) & Long (years; often multi-phase) \\
\hline
Political “deliverability” & Easy to showcase within one term; ideal for ward funds and patronage & Benefits arrive after the next election; harder to claim credit \\
\hline
Equity risk & High spatial inequality—favored streets improve first & Lower once built, but slow roll-out leaves many areas unserved for long periods \\
\hline
Key governance task & Prevent pork-barrel bias and ensure minimum standards everywhere & Create durable, cross-agency coalitions and financing mechanisms to finish the network \\
\hline
\end{tabularx}
\end{table}

For simple, readily targetable services—including intermittent water standposts, streetlights, or periodic waste removal—both conditional autonomy and collusion can facilitate some form of delivery. Under collusion, these goods may be distributed as political favors (\textit{localized patronage}), with politicians exercising full discretion. Under conditional autonomy, the same services achieve more even and reliable distribution, approaching \textit{universalized delivery}, though not achieving full universality (Figure~\ref{fig:universalism-model}).

However, for complex, networked infrastructure—such as piped water systems or sewerage networks—the distinction between conditional autonomy and collusion becomes critical. The technical demands of these goods eliminate the flexibility that makes targetable goods politically attractive. Conditional autonomy enables \textbf{precarious universalism}: broad, but not perfect, coverage that remains vulnerable yet substantially superior to the fragmented outcomes of collusive governance.

Without professional norms, strategic investments, and partial insulation from political pressures, complex infrastructure projects are prone to collapse. In collusive systems, non-targetable services experience \textbf{complete fragmentation}, precluding the possibility of stable, citywide solutions and reducing service delivery to patchy, politically motivated interventions rather than coherent development.

\begin{figure}[H]
    \centering
    \includegraphics[width=0.7\textwidth]{universalism_model.png} % <-- Replace with your figure
    \caption{Impact of Conditional Autonomy vs Collusion on Public Goods Delivery}
    \label{fig:universalism-model}
\end{figure}


Politicians operating within patronage systems actively resist non-targetable goods precisely because their technical characteristics eliminate the discretion essential for patronage politics. These goods cannot be selectively directed to reward supporters due to their inherent network properties. Unlike targetable goods, which allow politicians to choose which ward receives road improvements or better waste collection, non-targetable goods operate on an all-or-nothing basis, removing political discretion.

Under collusion, these characteristics lead to \textbf{complete fragmentation}. Politicians perceive no electoral advantage in championing non-targetable goods they cannot selectively distribute, while instrumentalized bureaucrats lack both autonomy and incentive to pursue them. The result is disconnected water projects and patchwork private systems, shaped by political calculations rather than technical coherence.

Conversely, conditional autonomy enables \textbf{incremental universalism}—the significant though imperfect expansion of networked services despite patronage pressures.
The differential performance between governance modes reveals why meso-level bureaucratic autonomy proves indispensable for delivering non-targetable goods:

\begin{itemize}
    \item \textbf{Technical complexity demands professional management:} Politicians typically lack the expertise to design integrated systems. Only bureaucrats with engineering training can navigate the technical requirements of networked infrastructure.
    
    \item \textbf{Long time horizons require institutional continuity:} While politicians face electoral cycles, meso-level bureaucrats provide the institutional memory essential for multi-year projects.
    
    \item \textbf{System integrity requires protection from diversion:} Non-targetable goods fail when funds are misappropriated or redirected. Conditional autonomy creates barriers against political interference.
    
    \item \textbf{Low political discretion creates bureaucratic opportunity:} Politicians often overlook non-targetable goods precisely because they offer no discretion in distribution and limited electoral payoff, inadvertently creating space for bureaucratic initiative.
\end{itemize}

The implications are profound. Reform efforts in weak institutional environments should focus precisely where conditional autonomy yields maximum impact: non-targetable, networked infrastructure. Rather than pursuing comprehensive institutional transformation, strategic interventions that protect meso-level bureaucrats managing water, sewage, or transportation systems can produce disproportionate welfare gains.

This targeted approach acknowledges political realities while exploiting blind spots in the patronage system. Politicians may resist granting broad bureaucratic autonomy but often overlook technical domains with limited patronage value. By empowering meso-level bureaucrats in these specific areas, reformers can achieve significant improvements without directly confronting the political foundations of patronage.

\section{Conditional Autonomy: A Relational and Negotiated Equilibrium}

The empirical evidence necessitates a fundamental reconceptualization of governance dynamics. Conditional autonomy represents neither a transitional phase nor a hybrid anomaly, but rather a stable, self-reinforcing equilibrium operating within patronage-dominated contexts. This equilibrium possesses its own internal logic, performance profile, and sustainability mechanisms that demand theoretical elaboration.

To restate, this dissertation conceptualizes conditional autonomy as simultaneously relational, evolving, and negotiated between politicians and bureaucrats at the meso level. Its relational nature emerges from ongoing interactions—a dynamic balance of power and mutual expectations—rather than formal institutional rules alone. Its evolutionary character manifests in the capacity to strengthen or erode as political and performance dynamics shift. Its negotiated aspect operates not through singular formal agreements but through implicit understandings, tacit accords, and occasionally explicit pacts that delineate bureaucratic discretion. Essentially, conditional autonomy grants bureaucrats limited independent authority contingent upon specific conditions: delivering measurable results, avoiding open defiance of core political interests, and demonstrating value to political patrons. This equilibrium occupies a middle zone between two theoretical extremes. It exhibits neither the full Weberian autonomy of bureaucracies impervious to political interference nor the complete capture characteristic of politicized bureaucracies. Instead, it represents a stable arrangement where both bureaucrats and politicians find maintaining partial bureaucratic independence rational and mutually beneficial.

The equilibrium's stability derives from bidirectional accountability and mutual constraint-benefit trade-offs. Politicians in conditional autonomy regimes restrain from excessive interference in specific bureaucratic domains, recognizing the performance benefits of limited autonomy. Reciprocally, bureaucrats utilize their discretionary authority to generate outcomes reflecting positively on political leadership while remaining cognizant that autonomy remains contingent. This creates a self-enforcing dynamic: bureaucrats maintain performance to protect their autonomy, while politicians preserve bureaucratic discretion to secure political dividends through improved service delivery.

In game-theoretic terms, this represents a repeated interaction achieving cooperative equilibrium. Politicians' credible threat of reverting to patronage control disciplines bureaucratic behavior, while bureaucrats' demonstrated capacity to deliver results deters political interference. This equilibrium demonstrates fragility yet resilience—vulnerable to shocks such as leadership changes or major policy failures but persistent while both parties perceive the arrangement as superior to alternatives.

This conceptualization resonates with Evans's notion of "embedded autonomy" in developmental states, though operating in a different context. Where Evans described bureaucracies combining autonomous structures with dense societal ties, conditional autonomy combines professional integrity with embeddedness within patronage politics. The bureaucracy operates not behind insulating walls but maintains professional zones within patronage pressures.

Conditional autonomy also resolves what Geddes termed the "politician's dilemma"—the tension between exploiting state resources for immediate political gain versus building competent civil services yielding longer-term benefits. Neither extreme serves politicians' interests optimally: fully autonomous bureaucracies undermine short-term electoral needs, while completely politicized bureaucracies sacrifice developmental capacity.

The logic of conditional autonomy becomes evident when examining performance across varying task complexity levels (Figure 4).

\begin{figure}[H]
    \centering
    \includegraphics[width=0.8\textwidth]{figure4_placeholder.png}
    \caption{Governance Mode Effect on Complexity of Good}
    \label{fig:governance_complexity}
\end{figure}

Figure 4 presents three distinct performance curves:

\begin{itemize}
    \item The Patronage Mode (blue line) demonstrates high effectiveness for simple tasks but experiences a steep, linear decline as complexity increases. This reflects patronage systems' ability to distribute targetable benefits efficiently while failing catastrophically at complex, networked tasks requiring technical coordination.
    \item The Programmatic Mode (red line) exhibits relatively low effectiveness for simple tasks but maintains steady performance as complexity increases, showing a gradual decline only at higher complexity levels. This pattern reflects programmatic systems' emphasis on universal rules and procedures over discretionary distribution.
    \item Conditional Autonomy (orange/yellow line) begins with high effectiveness for simple tasks— nearly matching patronage performance—but declines more gradually than patronage as complexity increases. Critically, it maintains superior performance to both programmatic and patronage modes across the middle range of complexity, demonstrating its unique advantage for moderately complex tasks.
\end{itemize}

This performance pattern reveals why conditional autonomy represents a viable equilibrium: it combines the distributional efficiency of patronage for simple tasks with sufficient technical capacity for moderately complex challenges. While it may not match programmatic systems at extreme complexity, conditional autonomy achieves "good enough governance" precisely where most essential public services operate—in the middle range of complexity where neither pure patronage nor rigid programmatic approaches perform optimally.

The graph thus illuminates why conditional autonomy emerges and persists: it offers a pragmatic balance between political imperatives and developmental outcomes, producing sufficient benefits for both politicians and citizens to ensure self-sustainability in patronage-dominated contexts.

\subsection{How Meso-Level Bureaucrats Exercise Conditional Autonomy}

The theoretical construct of conditional autonomy finds its practical expression through four distinct yet interconnected mechanisms which though alluded earlier are given fuller attention here. These mechanisms—bureaucratic brokerage, strategic intermediation, intertemporal stabilization, and normative institutionalization—represent not merely organizational functions but fundamental transformations in how governance operates under constraints. Together, they explain how mid-level officials convert limited autonomy into disproportionate governance improvements.

\paragraph{Intermediaries within the State: The Strategic Position of Meso-level Bureaucrats}

A fundamental challenge in state capacity is coordinating activities across line departments. This isn't just abstract theory—it's the everyday reality of getting water departments to talk to sanitation departments, or education officials to coordinate with health services when school feeding programs are implemented. Traditional "silo mentality" plagues bureaucracies worldwide, with each department jealously guarding its turf, budget, and authority. Meso-level bureaucrats occupy a uniquely advantageous position as intermediaries within the state apparatus itself—inverting our usual focus on brokers and middlemen operating outside formal state structures.

Their location in the administrative hierarchy enables critical coordination that neither elites nor street-level bureaucrats can achieve: Horizontally, they bridge departmental divides without the baggage that hampers others. Unlike top brass whose interactions are limited by protocol, status consciousness, and turf protection, mid-level officials can pick up the phone and call their counterpart in another department without triggering a bureaucratic crisis. They don't need formal memoranda or committee approvals to initiate problem-solving conversations. These officials share a common bureaucratic language and professional ethos that transcends departmental boundaries, allowing them to cut through red tape when line departments need to work together.

Vertically, these officials function as crucial conduits between ground realities and policy circles. Street-level bureaucrats—the teachers, nurses, and agricultural extension workers—possess intimate knowledge of implementation challenges but lack the authority to push these insights upward. They're trapped in their daily routines, processing cases one at a time. Department secretaries and directors, meanwhile, remain removed from operational messiness, focused on high-level policy formulation and political navigation. Mid-level bureaucrats translate between these worlds—bringing field observations to headquarters while interpreting policy intentions to implementers.

A municipal commissioner might convey to the state government that a water policy isn't working because local pipes are too old, while simultaneously explaining to local engineers why compliance with new state standards matters. This dual intermediary function makes mid-level bureaucrats the connective tissue of the state machinery. Unlike the classic focus on external intermediaries who operate between state and society (the ward boss, the fixer, the NGO facilitator), these internal intermediaries work within the administrative architecture itself, reinforcing connections between departments and across hierarchical levels. When they exercise conditional autonomy, they can address coordination deficits that formal structures often fail to resolve, creating pragmatic solutions to bureaucratic fragmentation even amidst pervasive patronage pressures. This intermediary function can be formalized as a mechanism of bureaucratic brokerage, whereby meso-level officials leverage their unique structural position to bridge otherwise disconnected departments and hierarchical levels, facilitating information flows and resource coordination that formal systems often fail to achieve.

% (The remaining sections – Boundary Spanners, Memory, Stewardship – would follow in same format, verbatim. Let me know if you want the rest split or posted now.)

\paragraph{Boundary Spanners: The Interstitial Role of Mid-level Bureaucrats}

While mid-level bureaucrats function as internal intermediaries, they simultaneously operate as crucial boundary spanners at the state-society interface (Williams, 2002; Honig, 2018). This interstitial positioning constitutes their second vital contribution to state capacity in contexts of conditional autonomy.

The boundary-spanning function emerges from a practical organizational dilemma: elite bureaucrats must maintain professional distance from potentially compromising interactions. Senior officials cannot regularly engage with contractors, local power brokers, or politicians with particularistic agendas without risking their professional standing or ethical reputation. Yet these engagements remain necessary for implementing policies and programs.

Mid-level bureaucrats fill this interstitial space, engaging with external actors that higher officials cannot directly approach. This arrangement functions through a tacit division of labor. Senior bureaucrats delegate boundary-spanning responsibilities to their mid-level subordinates, who negotiate with builders, contractors, elected representatives, and other external stakeholders.

A municipal commissioner might direct a deputy to handle negotiations with a construction firm rather than conducting these discussions personally. Similarly, politicians utilize these officials as intermediaries when pursuing both programmatic and particularistic aims, recognizing their operational knowledge and implementation authority.

The interstitial position creates a dual potential: mid-level bureaucrats can leverage these relationships to enhance program effectiveness and resource mobilization, or they may be instrumentalized by political or commercial interests seeking preferential treatment. Their position is inherently ambiguous, simultaneously offering opportunities for administrative entrepreneurship and avenues for co-optation.

This intermediary function can be formalized as a mechanism of strategic intermediation, whereby mid-level bureaucrats maintain legitimate standing with multiple principals simultaneously—their bureaucratic superiors, political authorities, and external stakeholders—creating protected spaces for administrative discretion even within patronage-dominated environments.

\paragraph{Institutional Memory Carriers: The Temporal Continuity Function}

While mid-level bureaucrats serve critical spatial intermediation functions—bridging departments horizontally and hierarchical levels vertically—they simultaneously perform an essential temporal function as carriers of institutional memory (Pollitt, 2009; Schön-Quinlivan, 2011). This contributes a third vital dimension to state capacity in contexts of conditional autonomy.

This function manifests in public administration systems worldwide, where projects spanning multiple years or even decades persist despite rapid turnover at senior leadership levels. While high-level officials may rotate frequently, mid-level bureaucrats maintain project continuity through their accumulated knowledge. Their institutional memory encompasses both formal aspects (regulation details, procedural requirements) and tacit knowledge (knowing which files matter, where documents are stored, how budgetary allocations work in practice).

This memory function addresses a fundamental governance challenge in developing contexts: the high turnover of senior officials that creates institutional amnesia. In many bureaucracies, elite positions rotate frequently due to political transitions or internal promotion patterns. Street-level bureaucrats, while potentially stable in their roles, lack the systemic perspective to preserve institutional continuity. Mid-level officials bridge this temporal gap, having sufficient longevity in their positions to accumulate deep organizational knowledge while possessing enough authority to deploy this knowledge effectively.

As Weber (1978) observed in his analysis of bureaucratic expertise, effective administration depends not merely on formal rules but on the "craft knowledge" that develops through immersion in organizational processes. Bourdieu's (1990) concept of habitus further illuminates how officials embody institutional practices, developing an intuitive feel for navigating complex bureaucratic fields. Mid-level bureaucrats exemplify this embodied knowledge, having internalized both explicit procedures and implicit norms that guide administrative action.

This intermediary function can be formalized as a mechanism of intertemporal stabilization, whereby mid-level bureaucrats maintain critical organizational knowledge across leadership transitions, ensuring policy continuity despite the discontinuous nature of political and elite administrative tenures in developing contexts.

\paragraph{The Professional Stewardship Function}

Beyond their intermediary roles, meso-level bureaucrats perform a fourth critical function through their cultivation of professional stewardship within state institutions. This dimension addresses a fundamental gap in traditional administrative theory that has predominantly characterized bureaucrats through rational choice frameworks—as utility-maximizing agents driven primarily by material self-interest or career advancement.

While strategic self-interest undoubtedly shapes bureaucratic behavior, field research reveals a more complex motivational landscape among mid-level officials. Their extended tenure within specific departments—often spanning decades in a single institutional setting—creates a distinctive relationship to their organizational context that transcends purely instrumental calculations.

Unlike elite bureaucrats who frequently rotate between postings or street-level workers with limited authority to shape institutional direction, mid-level officials develop profound identification with their departments. As one district officer articulated during field interviews: "I will spend my entire career in this department. If I don't improve it now, I'll have to live with the consequences for the next 30 years." This temporal horizon fundamentally alters their decision calculus, transforming short-term challenges into long-term investments in institutional effectiveness.

This professional stewardship emerges most strongly under conditions of partial autonomy, where bureaucrats retain enough discretion to experience meaningful agency while facing resource constraints that necessitate creative problem-solving. They cultivate what DiMaggio and Powell (1983) describe as "normative isomorphism"—professional standards that become internalized and reproduced across organizational contexts, creating resilience against patronage pressures.

This function can be formalized as a mechanism of normative institutionalization, whereby mid-level bureaucrats with extended departmental tenures develop and sustain professional values that persist beyond individual officeholders, creating institutional resilience that enables cumulative improvements in state capacity even within imperfect governance environments.


These four mechanisms—bureaucratic brokerage, strategic intermediation, intertemporal stabilization, and normative institutionalization—collectively illuminate how mid-level bureaucrats operate as the connective tissue of state capacity. By bridging structural, relational, temporal, and normative dimensions of governance, they create integration where fragmentation would otherwise prevail.

The combined effect is transformative: rather than bureaucracy that passively follows orders or actively serves private interests, we observe proactive, self-regulating innovation within political boundaries. The conditionality of autonomy ensures responsiveness, while the autonomy enables application of skill, knowledge, and coordination capacity to public problems.

This framework reconceptualizes the "missing middle" not as a theoretical gap but as a pivotal arena where state capacity is actively constructed through distributed agency. Mid-level officials navigate between formal mandates and informal realities, establishing horizontal connections, spanning external boundaries, maintaining institutional memory, and cultivating professional norms—practical statecraft that conventional theories overlook.

By illuminating these mechanisms, we gain both analytical leverage for understanding subnational variation and practical insights into fostering incremental improvements in state capacity, even within challenging institutional environments.

\section{A Model of Bureaucratic Authority in Weak Local States: The Divergent Paths of Governance}

To formalize the theoretical relationships outlined above, we develop a comprehensive model of bureaucratic authority pathways in weak local states, demonstrating how initial conditions and path dependencies produce divergent governance outcomes.

\begin{figure}[H]
    \centering
    \includegraphics[width=0.85\textwidth]{figure5_model.png} % Replace with actual filename
    \caption{How does conditional autonomy operate}
    \label{fig:figure5_model}
\end{figure}

Figure 5 presents a comprehensive model of bureaucratic authority pathways in weak local states, illustrating how initial political choices cascade into distinct governance equilibria. The model begins with a fundamental precondition: the existence of a weak local state characterized by limited institutional capacity and strong patronage networks.

The critical juncture occurs when political elites face pressure for reforms, either from citizens demanding better services or external actors promoting governance improvements. This pressure triggers a bidirectional decision point: whether political elites will grant autonomy to meso-level bureaucrats. The model reveals two distinct pathways, each with self-reinforcing dynamics that lead to divergent long-term outcomes.

When political elites deny bureaucratic autonomy, the system defaults to bureaucratic collusion. In this pathway, bureaucrats become fully instrumentalized by political patrons, leading to poor service delivery characterized by both targetable goods distributed as patronage and fragmented, non-viable attempts at complex infrastructure. The outcome crystallizes as localized patronage, where public resources serve narrow political interests rather than broad developmental goals.

When political elites opt to grant conditional autonomy, however, a more complex pathway emerges.

This choice triggers a cascade of positive developments: bureaucrats gain space to coordinate activities, build technical capacity, and address complex challenges that pure patronage systems cannot manage. The resulting improved service delivery—particularly in non-targetable, networked goods—generates a crucial feedback loop. Citizens and external actors provide political support for these improvements, which reinforces the political rationality of maintaining bureaucratic autonomy.

The model's key insight lies in the self-reinforcing nature of both pathways. Under collusion, poor performance legitimizes further political interference, creating a downward spiral. Under conditional autonomy, improved performance generates political rewards that sustain bureaucratic discretion, creating a virtuous cycle. This explains why jurisdictions with identical formal institutions can achieve dramatically different outcomes: initial political choices about bureaucratic autonomy set in motion path-dependent processes that become increasingly difficult to reverse.

The model also illuminates why conditional autonomy proves particularly crucial for delivering non-targetable goods. Only through sustained technical coordination and professional discretion can bureaucrats overcome the inherent challenges of networked infrastructure—challenges that patronage systems systematically fail to address. This theoretical framework thus provides both explanatory power for observed variations and practical guidance for reform efforts in weak institutional environments.

\subsection*{Summary}

This section accomplished two critical objectives. First, it identified significant blind spots in existing scholarship regarding variation in service delivery within patronage settings. The literature's focus on national-level institutions, abstract structural factors, and binary conceptions of state capacity has systematically overlooked the crucial role of meso-level bureaucratic dynamics in explaining subnational governance variation.

Second, it established a new theoretical framework to resolve this puzzle. The core argument posits that low-capacity, patronage-dominated polities perform unevenly because governance outcomes depend fundamentally on the interaction between political incentives and bureaucratic autonomy. When meso-level officials retain conditional autonomy—limited but real discretion within patronage systems—they can mobilize professional norms and inter-departmental authority to coordinate complex, citywide networks (water, sewerage, power) that individual politicians cannot parcel out as favors. This bureaucratic coordination improves service coverage and reliability, creating positive feedback loops that gradually institutionalize higher-capacity governance.

Conversely, where such autonomy is absent and officials operate in fully collusive relationships with politicians, service delivery fragments into localized patronage enclaves, trapping the state in low-performance equilibria.

The section also provided measurement strategies, a causal model, and specified the characteristics of meso-level officers and conditions under which they can innovate.

The next section transitions to examining the empirical context from which these theoretical insights emerge—Indian cities. It first establishes what Indian cities represent as cases for theory development, demonstrating their value for state capacity research and comparative politics more broadly. It then presents detailed case studies of Bhopal and Lucknow, providing empirical validation for the theoretical model developed here.

The next section moves to understanding the context from which these theory has been driven – Indian cities. First, it provides you with an overview asking – what is Indian city a case of? And how can it be a site for theory generation for state capacity in particular and comparative politics in general. It then discusses the cases I study, Bhopal and Lucknow – showing the validity or proof of concept of this model.

\part{Part B}

\section{What is the Indian City a Case of: Indian Cities and the Paradox of State Capacity}

\subsection{Fragmented and Uneven State Capacity in Urban India}

Indian cities powerfully illustrate the fragmented and uneven nature of state capacity in the Global South. In development theory, scholars often note that state authority does not extend uniformly across territory or social groups. Guillermo O'Donnell's notion of "brown areas," for instance, pointed to enclaves within democracies where the rule of law and public institutions are weak.

This spatial and institutional fragmentation creates precisely the conditions where conditional autonomy emerges as a governance mode—mid-level bureaucrats must navigate multiple, overlapping authorities through negotiated practices rather than clear lines of command. This empirical reality makes Indian cities ideal cases for examining how the meso-level bureaucrats theorized in Section 1 navigate governance challenges. What conventional theoretical lenses dismiss as dysfunction instead creates the interstitial spaces that make conditional autonomy not just possible but necessary for effective governance.

Indian urban centers exemplify this pattern of institutional fragmentation in particularly stark terms. Within a single metropolis, pockets of high-capacity governance coexist with areas where the state's reach is tenuous and public services are deficient. This unevenness is not merely a rural-urban divide, but an intra-urban reality—a city of enclaves and exclusions rather than a monolithic unit of governance. In short, Indian cities encapsulate a broader paradox: a state that is impressively present in some domains and absent in others, challenging conventional models that treat state capacity as a homogenous national trait.

Crucially, this fragmented state capacity is not a static accident of history—it is produced and reproduced by governance structures. Indian cities operate within a federal system where local governments have formally defined powers, yet practically constrained autonomy. The result is often "islands" of capacity—effective delivery of certain services or in certain locales—surrounded by vacuums or overlaps elsewhere. Such outcomes reflect more than just uneven investment; they point to how the machinery of the state itself is fragmented, with multiple agencies and levels of government struggling (or failing) to coordinate. Indian cities thus personify the idea that state capacity is not monolithic—it varies widely within countries and even within cities, defying any singular characterization of "strong" or "weak" state. This reality provides fertile ground to test and refine theories of state capacity and development that have often been drawn from more unified national experiences.



\subsection{Overlapping Governance and Public Goods Delivery}

A key reason for the patchwork performance of Indian cities is their complex governance structure – a maze of overlapping jurisdictions and divided authority. Unlike cities in many other democracies, Indian municipalities do not enjoy exclusive control over urban affairs. Instead, critical functions are shared (and sometimes contested) among multiple institutions: city corporations, state government departments, public sector boards (for water, transport, housing, etc.), and even special development authorities. This multiplicity leads to ambiguous responsibility. Who is accountable for providing water, managing traffic, or maintaining drains in a city? Often it is not one entity but several, each with a piece of the task. The consequence is duplication in some areas and vacuum in others – a recipe for inefficient and unequal public goods delivery. Four structural features characterize urban governance in India and shape public goods provision:


\subsubsection*{Multiplicity of Agencies}

Indian cities operate under fragmented governance structures, with state- created parastatals functioning parallel to elected municipal bodies. These specialized agencies—water boards, development authorities, slum clearance boards—report directly to state governments, bypassing local democratic oversight. As illustrated in Figures 1-2 this fragmentation manifests starkly in spatial planning arrangements across major metropolitan regions. In Mumbai, 13 sub-regional planning authorities operate alongside the Municipal Corporation. Hyderabad's planning authority absorbed four erstwhile agencies yet still operates eight plans prepared by multiple authorities which remain in force. In Bangalore this number is 8 and Chennai even with a relatively simpler structure, divides planning between the development authorities master plan and a staggering 57 high-growth areas requiring detailed development plans (Table 3).

What appears as an organizational arrangement is fundamentally a political problem of sovereignty. These parastatals were ostensibly created to enhance technical capacity, but in practice have "usurped functions and revenues that ought to have been the domain of cities," undermining what Heller calls the "binding of the state" to citizens through participatory processes. Most critically, control over urban planning and land use—the foundation of meaningful urban governance—remains centralized in 80\% of states/UTs. This institutional arrangement creates precisely the governance vacuum where conditional autonomy becomes significant: mid-level bureaucrats must navigate between formal constraints and local realities, mediating between state directives and community needs.

Without territorial authority over land, municipal governments cannot function as strategic actors, instead becoming spaces where bureaucrats exercise conditional discretion within narrowly defined parameters—often creating the very governance gaps that require bureaucratic entrepreneurship to address. This institutional fragmentation directly illustrates the governance context theorized in Section 2, where meso-level bureaucrats must employ bureaucratic brokerage to coordinate across siloed agencies and departments. Together, these features create a governance environment where mid-level officials must navigate both horizontal fragmentation across agencies and vertical disconnection from democratic accountability.


\begin{figure}[H]
    \centering
    \includegraphics[width=0.8\textwidth]{figures/mumbai_planning.png}
    \caption{Multiple Spatial Planning Arrangements in Mumbai Metropolitan Region. The map shows 13 sub-regional planning authorities operating alongside the Municipal Corporation, illustrating the extreme fragmentation of governance. Source: Regional Plan for Mumbai Metropolitan Region (1996--2011), MMRDA.}
    \label{fig:mumbai_planning}
\end{figure}

\begin{figure}[H]
    \centering
    \includegraphics[width=0.8\textwidth]{figures/hyderabad_planning.png}
    \caption{Multiple Spatial Planning Arrangements in Hyderabad Metropolitan Region. Despite HMDA's formation by merging four erstwhile planning authorities, multiple planning regimes remain in force. Source: CPR Analysis.}
    \label{fig:hyderabad_planning}
\end{figure}


\subsubsection*{Administrative Dominance: When Bureaucracy Eclipses Democracy}

Urban governance across India exhibits a consistent pattern: turnout in municipal elections stays stubbornly low even as citizens energetically “work the state” through personal networks, political brokers, and direct bureaucratic appeals. In this context, appointed officials regularly eclipse elected representatives, and the ubiquitous commissioner system has become less a neutral administrative device than a political arrangement in which bureaucratic hierarchies substitute for democratic channels. Municipal Commissioners—typically drawn from the elite Indian Administrative Service (IAS)— control staff appointments, budgets, and day-to-day service delivery with only minimal local oversight. Yet my calculations show that these IAS commissioners serve in any given city for a median tenure of
 


just 16 months, whereas mid-career state-cadre (“meso-level”) officers in Madhya Pradesh (MP) and Uttar Pradesh (UP) remain in comparable municipal posts far longer (18 months for UP and 24 months for MP). It is precisely these longer-serving, locally embedded administrators who form the empirical core of this study: their extended tenures furnish the organizational memory and discretionary space that make—or break—urban governance outcomes (Figure 6).

Meanwhile, elected Mayors serve largely ceremonial functions despite being the formal face of city government. This creates what might be called a ``sovereignty gap'' in urban governance. While formal democratic processes exist—elections are held, councils meet, debates occur—the substantive authority to implement decisions flows through bureaucratic rather than democratic channels. The constitutional amendment of 1992 that formally recognized urban local bodies (city governments) attempted to bridge this gap, but its implementation paradoxically reinforced administrative dominance by giving state governments discretion over the actual powers devolved to cities.

The governance arrangement where bureaucratic hierarchies—at elite, meso, and frontline levels— exercise greater authority than elected bodies creates precisely the conditions this dissertation examines —neither fully autonomous (as they must respond to state-level political pressures) nor completely captured (as they retain significant discretionary authority over implementation). Unlike classic governance models where bureaucrats simply implement political decisions, officials at all levels must actively interpret, adapt, and negotiate formal directives to accomplish basic functions. This distributed bureaucratic authority, particularly at the meso-level, becomes critical for understanding how cities actually function despite their institutional constraints.

The consequences run deeper than institutional design. This governance model shapes the very nature of urban politics. Elected councilors, unable to directly control administration, evolve into specialized intermediaries who help constituents navigate bureaucratic processes rather than crafting policy. The bureaucracy, nominally apolitical, becomes the de facto center of political dealmaking and resource allocation. This inversion creates governance spaces where informal negotiations between bureaucrats and special interests often matter more than formal democratic deliberations.

The result is a governance ecosystem where democratic energy exists but lacks institutional channels for effective expression. Urban residents experience government primarily through bureaucratic encounters rather than democratic representation. This arrangement persistently undermines the development of vibrant local democracy while simultaneously requiring bureaucratic innovation to overcome the coordination failures it creates—a self-reinforcing cycle that maintains administrative dominance while necessitating the very discretionary bureaucratic practices it supposedly constrains.

\section{Fiscal Emasculation: The Financial Incapacitation of Urban Governance}

City governments in India face crippling fiscal constraints that fundamentally undermine their governance capacity. This financial powerlessness represents not merely a technical deficiency but a structural governance arrangement where municipalities possess formal responsibilities without corresponding financial tools—what scholars call ``unfunded mandates'' (Ahluwalia, 2019).

The fiscal marginalization is quantitatively severe: municipal governments control only 1-2\% of public expenditure, compared to 20-35\% in other major federal democracies (Mathur, 2021). Figure 7 demonstrates this stark disparity—while local governments in countries like the United States, Spain, and Brazil control between 20-45\% of government expenditure, India's local governments manage a mere 3\%, the lowest among all compared nations. This constraint operates through multiple institutional mechanisms. This constraint operates through multiple institutional mechanisms: constitutional provisions for predictable resource transfers remain unimplemented by most states; taxation authority is severely curtailed; and borrowing requires explicit state approval, effectively foreclosing municipal access to capital markets.

\begin{figure}[htbp]
\centering
\caption{Share of Local Government Expenditure Across Countries}
\label{fig:fiscal_constraint}
% Insert your figure here
\end{figure}

Three consequences directly flow from this fiscal architecture. First, cities depend on discretionary, politically-mediated transfers rather than stable, rules-based funding. Second, financial management deteriorates as municipal corporations routinely miss budget targets by 40-60\%, creating a governance environment where planning becomes performative rather than operational. Third, accountability mechanisms collapse as responsibility and authority disconnect—elected officials face citizen demands they lack resources to address.

Their resulting practices—informal resource mobilization, negotiated rule enforcement, contingent service delivery—represent not corruption but strategic adaptation to systemic constraints. The comparative perspective reveals this is not inevitable: South African municipalities receive constitutionally guaranteed revenue shares, while Brazilian cities control significant local tax bases (Shah \& Shah, 2020).

Without addressing this deliberate fiscal incapacitation, the formal trappings of urban democracy— elections, councils, citizen participation—will continue to mask the substantive reality: cities that cannot govern effectively because they cannot finance even their most basic mandates. This financial powerlessness exists within and reinforces a broader governance framework that systematically concentrates resources at state and national levels, leaving cities perpetually under-resourced despite their growing economic importance and population.

\section{Residual Urbanism: Rural Primacy \& The Systematic Marginalization of Urban India}

India's governance exhibits a profound rural bias that manifests in three distinct ways. First, restrictive classification criteria systematically undercount urban populations. While official statistics place urbanization at 31\%, my earlier work has shown that applying internationally comparable definitions would reveal that India is approximately 40\% urban (Joshi \& Pradhan, 2018) The peculiar category of ``census towns''—settlements that meet demographic criteria for urban status yet remain administratively rural—exemplifies this statistical distortion\footnote{India employs one of the most restrictive urban definitions globally. Urban areas are divided into Statutory Towns (with elected urban local bodies) and Census Towns (meeting three criteria: population $\geq$5,000, density $\geq$400 persons/km², and $\geq$75\% male workforce in non-farm sector). Despite qualifying as urban, Census Towns remain administratively rural, governed by rural local bodies.}

Second, institutional incentives drive deliberate misclassification. State governments strategically maintain rural classifications because rural status confers greater financial resources with fewer administrative constraints. The most striking example occurred in 2004 when Tamil Nadu—then India's most urbanized state—reclassified over five hundred urban areas as rural to obtain more funds from the Union government (Sivaramakrishnan, 2013). This represents not statistical error but rational response to a governance framework that privileges rural development.

Third, and most fundamentally, Indian policy conceives urban space as merely residual defined not by what it is but by what it is not (rural). Unlike the ``urban bias'' literature pioneered by Michael Lipton (1977) that positions cities as privileged in developing countries, India's reality presents the inverse: a governance system that systematically favors rural interests while neglecting spaces where economic and demographic growth concentrate most intensely. While Varshney's (1993) critique demonstrated that Lipton's original formulation neglected key mediating factors—particularly democratic institutions providing rural populations electoral leverage, technological changes like the Green Revolution enhancing rural political power, and identity politics fragmenting simple rural-urban coalitions—the Indian case nonetheless represents a clear inversion of the urban bias thesis.\footnote{Varshney's analysis emphasizes that democratic institutions fundamentally alter the political landscape by giving rural populations electoral leverage absent in Lipton's original framework. The nature of political regimes, technological transformations, and identity-based mobilization all mediate the expression of sectoral interests in ways that complicate any simple model of structural bias.}

This rural primacy creates profound governance consequences that extend beyond simple resource allocation. Most critically, it generates a systemic urban representation deficit in legislative bodies. Due to the prolonged freeze on constituency delimitation based on 1971 Census data, urban areas now suffer severe under-representation. Bengaluru, despite accounting for 14\% of Karnataka's population, holds only 10-12\% of legislative seats; Pune's urban areas receive only 2.7\% of Maharashtra's assembly seats; and nationally, only 15\% of Lok Sabha constituencies can be classified as majority urban despite approximately 40\% of the population living in urban areas.\footnote{This representation deficit violates the democratic principle of `one person, one vote, one value' due to vastly different constituency sizes. The freeze on seat reallocation after 1971, extended through successive constitutional amendments, creates a political imbalance where urban legislators represent significantly larger populations than their rural counterparts. Combined with weak fiscal autonomy and heavy dependence on intergovernmental transfers, this political under-representation exacerbates urban governance challenges.} This structural under-representation compounds fiscal weaknesses, as Urban Local Bodies depend on central transfers rather than robust own-source revenues, creating gaps in infrastructure and service delivery precisely where population pressure is most acute. The result is a governance ecosystem that systematically undermines urban capacity while simultaneously requiring adaptive strategies from diverse actors—bureaucrats, elected officials, civil society organizations, and citizens—who must navigate this marginalization through the complex mechanisms of negotiation and intermediation this dissertation examines.

These four structural features collectively create what might be termed an `accountability vacuum' in urban governance. The principal-agent relationship fundamental to democratic accountability—where voters (principals) hold elected officials (agents) responsible for service delivery—becomes fundamentally distorted. Municipal councilors lack authority over bureaucrats who report to state capitals; service-providing agencies answer to multiple masters; and financing flows through discretionary rather than rules-based channels, leaving elected representatives unable to fulfill their democratic mandates.

This accountability fragmentation creates precisely the governance space where conditional autonomy emerges as a necessary adaptation. Within these institutional constraints, mid-level bureaucrats must deploy the four mechanisms theorized in Section 2: bureaucratic brokerage to coordinate across fragmented agencies, strategic intermediation to interpret ambiguous mandates, intertemporal stabilization to maintain policy continuity amid political flux, and normative institutionalization to establish informal rules that compensate for formal gaps. Their effectiveness depends not on following rigid hierarchies but on navigating relationships with multiple stakeholders—state officials, elected representatives, community organizations, and private interests.

The resulting governance model explains both the persistent failures and unexpected successes observed in Indian cities. Services deteriorate where bureaucrats cannot bridge these accountability gaps, creating neighborhoods that fall through institutional cracks. Conversely, effective service delivery emerges where entrepreneurial officials successfully navigate fragmented structures through conditional discretion—creating `islands' of capacity within otherwise dysfunctional systems. Understanding these dynamics requires moving beyond institutional design to examine how bureaucratic agency operates within and around formal constraints, precisely the contribution this dissertation makes to theories of state capacity in constrained environments.

\section{Bureaucratic Autonomy versus Political Accountability}

The governance configuration of Indian cities raises a classic question in development theory: how does one strike the right balance between bureaucratic autonomy and political accountability, particularly in contexts of severe fiscal constraints, fragmented legal authority, and financial dependence. The Indian case suggests that this balance is especially elusive in urban contexts, and it complicates prevailing models like the principal--agent framework or Peter Evans's notion of ``embedded autonomy.''

Bureaucratic autonomy—the insulation of civil servants from politicization—is traditionally seen as a virtue of capable states. In Indian cities, however, we often see the downsides of excessive autonomy and of too little autonomy, coexisting uneasily. The autonomy of city bureaucrats in India exemplifies our conditional autonomy framework from Section 2. These officials operate with partial independence from local electoral pressures, shielded by their appointment through state/national civil services. This creates precisely the type of negotiated equilibrium we theorized, where bureaucrats maintain professional discretion within politically defined boundaries, neither fully insulated nor completely captured. This could prevent short-termism or clientelist distribution at the local level. Yet, this autonomy from local accountability can also morph into detachment and unresponsiveness. A bureaucrat who rotates into a city for a couple of years might feel answerable primarily to superiors in the state capital rather than to the city's populace. If the Commissioner prioritizes state government directives or personal career considerations (e.g. avoiding controversy), local needs may be sidelined. In this sense, bureaucratic autonomy in Indian cities often means autonomy from the public, which can undermine the very purpose of public service.

On the other hand, bureaucrats in Indian cities are not wholly autonomous in a Weberian sense – their functioning is deeply enmeshed in political currents, only the currents come from above and outside the city. State-level politicians (chief ministers, ministers, legislators) routinely intervene in city matters, including in personnel decisions. Transfers of municipal officers at the behest of politicians are commonplace, sending a clear signal that stepping on powerful toes can be career-ending. There are numerous instances of municipal commissioners or ward engineers being shunted out because their actions (like removing illegal encroachments or enforcing rules impartially) displeased an influential politician or interest group. This dynamic fits a principal-agent logic, but with a twist: the bureaucrat becomes an agent to multiple principals – not just the lawful political principal (the city's elected body) but also to higher-level political actors. The result is a tenuous bureaucracy, adept at surviving through accommodation. The famed ``steel frame'' of the All-India Services in practice bends to political winds in the urban context. Autonomy, in the Evans sense, is compromised when bureaucrats must anticipate political reactions to every move. But unlike a clean principal-agent hierarchy, here the lines of authority are blurred: city bureaucrats navigate a dual accountability – upward to state leaders, and nominally downward to local citizens, often pleasing the former at the expense of the latter.

This scenario complicates principal--agent models. Classic theory envisions elected officials (principals) setting goals and monitoring bureaucrats (agents) to ensure implementation aligns with voters' interests. In Indian cities, who is the principal? The local councillor has formal authority but limited power; the state minister has actual clout but is not directly accountable to the city's voters for municipal services. The agent (bureaucracy) thus juggles instructions from multiple principals with differing priorities. Moreover, the agents themselves have some autonomy to pursue their own agendas or biases, especially given the opacity of who is in charge. Such a system can breed accountability confusion and even perverse incentives – for example, a municipal officer might delay action on a local issue until a state politician intervenes, thus securing favor with the latter. The principal-agent framework can be stretched to accommodate multiple principals, but Indian cities suggest we might need more networked or polycentric models of governance to describe reality. Authority is dispersed and negotiated rather than neatly delegated.

Another theoretical lens is provided by street-level bureaucracy, Michael Lipsky's idea that front-line public servants effectively make policy through their discretionary actions. Indian cities present textbook conditions for street-level bureaucracy to flourish. This discretionary power operates within a context of severe human resource deficits. Indian municipalities are dramatically understaffed compared to international counterparts—Bengaluru employs only 317 staff per 100,000 residents and Delhi about 601, compared to nearly 3,000 in London and 6,000 in New York. As Figure 8 illustrates, India's local governments have the lowest share of public employees among comparable emerging economies, with state governments absorbing the vast majority of administrative capacity. This staffing crisis is exacerbated by vacancy rates often exceeding 30-40\%, alongside significant skill gaps in specialized areas like urban planning, transport engineering, and financial management. With limited professional training and state governments controlling key appointments, municipal bureaucrats operate in an environment of both scarcity and constraint. Within this capacity vacuum, the few available front-line workers (clerks, engineers, inspectors, police constables) end up exercising enormous discretion in allocating public goods. Because high-level policies and directives often do not align or are infeasible in the face of resource constraints, street-level officials become de facto policymakers.

\begin{figure}[htbp]
\centering
\caption{Distribution of Public Employees Across Government Levels}
\label{fig:employees_distribution}
% Insert your figure here
\end{figure}

For example, a ration shop officer in a city may decide which families receive scarce subsidized foodgrains, or a junior engineer may informally prioritize repairing the water line in one neighborhood over another. In Mumbai's water distribution system, valve operators (chabiwalas) exemplify this phenomenon with particular clarity—these operators literally hold the keys to the city's water supply, determining through their daily valve adjustments which neighborhoods receive water, when they receive it, and with what pressure (Anand, 2017). Their decisions transform physical infrastructure into a mechanism of social stratification, creating what Anand terms ``hydraulic citizenship''—a form of belonging defined by one's ability to secure reliable water access. These ground-level decisions reflect not only official guidelines but also personal judgments, prejudices, and pressures. In a context where formal accountability is weak, street-level officials are susceptible to influence by local power brokers or their own rent-seeking interests. The phenomenon of ``speed money'' – small bribes to get things done – exemplifies how street-level bureaucrats can gatekeep services, effectively ruling in practice over who gets what.

Thus, while Indian cities suffer from high-level fragmentation, it is often the street-level functionaries who stitch together some semblance of order, albeit through improvised, ad-hoc adjustments rather than systematic policy. This confirms Lipsky's insight that the real policies delivered to citizens are often decided on the street, not in official plans – but it also shows how such discretion, in the absence of accountability, can perpetuate uneven service and even arbitrariness.

In sum, the Indian urban experience pushes us to refine ideas of bureaucratic autonomy and principal-agent relations. It suggests that autonomy must be paired with embeddedness in accountability structures – a point that recent research underscores. For instance, a study of Indian districts found that bureaucrats who hail from the area they administer (hence more embedded locally) can improve public goods provision, but only when local accountability (media, literacy, civic engagement) is strong; otherwise, local ties may simply enable capture by elites (Bhavnani and Li 2018). This resonates in cities: when municipal bureaucrats work in isolation from citizen voice (no strong accountability mechanisms like transparent feedback or empowered councillors), their autonomy veers into unresponsiveness; yet if they were fully subject to local politics without institutional safeguards, service delivery could devolve into parochial patronage. Indian cities thus exemplify the challenges of achieving balanced bureaucratic autonomy – that is, independence from partisan exploitation without detachment from public accountability.

\section{The Role of Clientelism and Informal Mediation}

Given the structural constraints and the often-impersonal stance of the formal state apparatus, how do ordinary urban citizens in India secure the public goods and services they need? The answer, in many cases, lies in informal networks, patronage, and clientelism filling the breach. When rules and institutions do not guarantee basic necessities, people turn to what Partha Chatterjee famously described as ``political society'' – the realm where citizens mobilize not just as rights-bearing individuals, but as members of patron-client networks, trading political support for material benefits.

Operating within this context of limited formal power and a politically fragmented urban society— where caste, religion, and class significantly shape political mobilization and voter allegiances— councilors adopt patronage politics as a rational strategy for political survival. Unable to effect broad policy changes, they leverage their limited, informal influence over the distribution of targetable non-networked goods like fresh tar on a patch of road, trash management, or lamp posts controlled at the local level. Within this constrained governance environment, a complex ecosystem of brokers—middlemen, touts, and fixers—has embedded itself into public service delivery systems. As recent research demonstrates, these brokers act as gatekeepers who ``block or expedite access to public services based on payment of a fee'' in cities like Delhi.

This reframing of clientelism as rational adaptation rather than mere corruption aligns with significant theoretical developments in political science and institutional analysis. Recent scholarship by Jennifer Gandhi, Adam Przeworski, and Beatriz Magaloni demonstrates how seemingly clientelistic practices often emerge as rational responses to specific institutional constraints rather than from corrupt motivations. Similarly, Rebecca Weitz-Shapiro's research on Argentina shows that clientelism persists precisely because it solves concrete governance problems in environments where formal mechanisms fail. Urban data from India for 14 cities substantiates this perspective: corporators receiving predominantly positive evaluations (60\% viewed as serving all constituents) suggests citizens themselves distinguish between corrupt rent-seeking and legitimate, if constrained, service provision.\footnote{Source: Authors' calculation based on Citizenship, Inequality, and Urban Governance (CIUG) Project, Brown University} Importantly, this ``constituency service clientelism'' operating through elected representatives differs significantly from more problematic forms involving illegal kickbacks or vote buying. By recognizing clientelism as a rational institutional adaptation—albeit one with significant governance costs—we gain deeper insight into why these practices persist despite decades of reform efforts. The persistence of such arrangements reflects not primarily moral failings but structural constraints that make alternative governance modes practically infeasible for both representatives and citizens.

Our survey data across 14 Indian cities provides empirical evidence that refines our understanding of how this system operates in practice. Contrary to theoretical expectations that emphasize state-level legislators or informal middlemen, citizens overwhelmingly identify municipal corporators as their primary channel for accessing public services. As Figure 9 demonstrates, in smaller cities, corporators remain the primary channel for accessing services with 51\% of residents identifying them as their first point of contact, while in large metropolitan areas this drops to 25\%. Even more surprising, citizens maintain largely positive views of these relationships—60\% believe corporators serve all constituents equitably rather than particular communities (13\%) or themselves (18\%). This suggests that what scholars often frame as clientelism is experienced by many citizens as legitimate constituency service.

\begin{figure}[htbp]
\centering
\caption{Who do you think is most important to deliver services to you?}
\label{fig:service_delivery}
% Insert your figure here
\end{figure}

Indian cities are frequently portrayed as arenas of vibrant electoral competition, yet this narrative masks two important realities. First, voter turnout in municipal elections consistently lags behind state and national contests. As Figure Y illustrates, municipal elections average only 49\% turnout compared to 62\% in parliamentary and 59\% in state elections, reflecting what Dasgupta (2015) terms ``the local democracy gap.'' Second, beneath the formal electoral processes lies a hustle for direct access: access to water, housing, electricity, sanitation, jobs, and permits through personal intervention rather than systemic improvements.

\begin{figure}[htbp]
\centering
\caption{The Local Democracy Gap: Average Voter Turnout by Election Type}
\label{fig:voter_turnout}
% Insert your figure here
\end{figure}

These intermediation practices operate within an institutional environment that severely constrains corporators' formal authority. Shadowing corporators in Lucknow revealed the limitations of their role. As one candidly expressed, ``We are firefighters, trash pickers, and drain openers. That is our role (hum aag bhujane wale hai, kuda uthane wale hai, naali khulwane wale hai).''\footnote{Quote from a city councilor in Lucknow, 2022} This statement vividly illustrates how local politicians are forced to assume administrative functions typically handled by bureaucrats. In the absence of substantive policy-making authority, corporators rely on patronage not primarily as a corruption strategy but as their only viable means of exercising influence. The structural constraints within which they operate—approximately 31\% of India's 4,766 municipal governments were effectively inactive as of 2023—transform even formal representatives into brokers within a dysfunctional system\footnote{Authors calculation using Local Self Government Directory and Comptroller and Auditor Generals Report}.

These personalized interventions represent classic examples of clientelistic exchanges – the voter gets a problem solved as a favor, and the politician hopes to get loyalty (or at least credit) in return. For marginalized citizens who lack formal paperwork or find themselves ignored by the system, a patron who can bend the rules becomes a lifeline. Research by Auerbach and Thachil (2018) reveals that urban brokers, unlike their rural counterparts, provide personalized and highly targeted services such as dispute resolution and assistance with documentation. Clientelism, in this sense, functions as a coping mechanism that compensates for the state's uneven performance – a way to navigate a fragmented system by leveraging personal influence rather than universal rights.

The consequences of this system for effective urban governance are profound. It undermines strategic, city-wide planning and leads to inequitable and inefficient service delivery, frequently prioritizing politically connected groups or areas while neglecting systemic improvements. Accountability becomes diffused, shifting from formal performance standards to informal, often opaque, transactional relationships, which can obscure responsibility and potentially foster corruption. Meaningful citizen participation is often curtailed, reduced to making particularistic demands on councilors rather than engaging in broader policy debates or holding the ULB accountable for overall city performance.

However, the clientelistic model has profound limitations that become increasingly evident in India's big cities. First, it is inherently particularistic—only some people benefit, often at the expense of others. In an urban constituency of say 300,000 people, even the most energetic representative can directly assist only a fraction. As research shows, politicians inevitably end up solving problems for some constituents while alienating others, even among their core supporters.

Second, this system does not scale well as cities grow. As Simon Chauchard notes, the increasing size of urban electorates and the complexity of needs make quid-pro-quo arrangements hard to sustain – there are simply too many constituents and issues for patronage to cover, and urban voters are often more heterogeneous and less tied to traditional patron-client bonds than rural voters. This scaling problem is reflected in our survey data, which shows that positive perceptions of corporators decline significantly from small cities (72\%) to large metropolises (55\%). Modern cities thus dilute the effectiveness of old-style machine politics; while clientelism is still present, it struggles to cope with the breadth of urban service deficits.

The institutional dysfunction that necessitates these clientelistic relationships creates a suboptimal equilibrium where no party truly benefits. When formal governance mechanisms break down—councils are dissolved, scheduled meetings fail to occur, and administrative processes stall—citizens have no choice but to reach out to councillors as their first line of action. These councillors, in turn, become overwhelmed by individual requests they lack the institutional capacity to address systematically. The resulting arrangement traps both citizens and representatives in a perpetual cycle of ad-hoc problem-solving that diverts attention from the structural reforms needed for effective urban governance.

Additionally, clientelism and informal mediation can undermine institutional development. If everyone must go through a ``broker'' to get basic services, there is less public pressure to improve the system for all. In fact, politicians themselves can become invested in the status quo of dysfunction, since it validates their role as indispensable fixers. This dynamic is captured in the phrase: ``The rich have peers, the poor have patrons.'' Wealthier or more educated city dwellers can often navigate bureaucracy through formal or professional networks (or purchase private alternatives), whereas the poor depend on patronage. Such patterns entrench inequalities of access. They also blur the lines between the formal state and informal practices: many bureaucrats expect political recommendations or informal payoffs to expedite work, effectively entwining clientelism with everyday administration.

These findings complement and extend Chatterjee's conceptualization of political society. While Chatterjee correctly identifies how marginalized groups must negotiate with the state through political rather than purely rights-based channels, our data reveals the central role of formal elected representatives within these negotiations. Rather than operating primarily through community-based mobilization, urban residents frequently engage directly with corporators as individuals seeking services. This direct engagement with corporators, combined with their positive perception despite limited formal powers, suggests a more complex relationship between citizens and the state than is captured in standard clientelism frameworks.

At the same time, one should not ignore the adaptive ingenuity in these arrangements. Urban residents form associations, slum federations, or use local media to amplify their issues, essentially trying to collectivize bargaining with the state when possible, not just individualize it. There are instances where community-led efforts got the municipality to install street lights or toilets, effectively compelling the state to respond without traditional patronage. Such civic initiatives interact in complex ways with clientelism – sometimes they supplant it, other times they themselves seek the patronage of a sympathetic official.

In summary, informal governance – through clientelist ties, brokers, and on-the-ground improvisation – is a double-edged sword in Indian cities: it provides relief and navigation in a fragmented system, but it also reinforces fragmentation by working around systemic solutions. Any analysis of state capacity and public goods in these cities must therefore grapple with this symbiosis of formal structures and informal strategies, recognizing that what appears as patronage often represents the only viable channel through which citizens can access basic services in a governance system that structurally constrains democratic mechanisms.

\section{Fragmented but Functioning: Rethinking State Capacity in Indian Cities}

While the previous sections establish that Indian cities exhibit profound deficiencies by most conventional measures of ``state capacity''—fragmented authority, weak fiscal autonomy, administrative undermanning, and structural constraints—a puzzling reality persists: these same cities continue to function and even thrive in crucial domains. Basic services, though unevenly distributed, reach millions of residents; urban economies generate significant growth; and innovative policy approaches regularly emerge from this seemingly dysfunctional landscape. This paradox—governance systems that appear broken yet remain operational—raises a fundamental question: How do Indian urban systems sustain themselves amid such institutional disorder? The conventional metrics of state capacity, with their Weberian ideals of bureaucratic coherence and hierarchical clarity, fail to explain this resilience.

\subsection{The Empirical Paradox: Progress Against All Odds}

What makes this theoretical puzzle empirically compelling is that Indian cities have demonstrated measurable improvements in service delivery over decades despite their institutional constraints. National Sample Survey data since 1992 reveals steady gains in access to basic services even under stringently defined parameters: piped tap water within premises and safe sanitation (toilets connected to sewers, septic tanks, or ventilated pit latrines meeting UNICEF standards) (Figure 11). This empirical reality fundamentally challenges purely pessimistic narratives of governance failure that dominate much of urban studies literature.

\begin{figure}[htbp]
\centering
\caption{Public goods in Indian cities over three decades. Authors Calculation}
\label{fig:public_goods}
% Insert your figure here
\end{figure}

I argue that state capacity in the Indian urban context is not a static attribute that municipalities either possess or lack, but rather an emergent property of complex, adaptive systems. It manifests through ongoing processes of negotiation, contestation, and co-production among diverse actors working within and around institutional constraints. What formal institutional analyses might dismiss as dysfunction— overlapping jurisdictions, informal arrangements, and ad hoc problem-solving (jugaad)—often represents creative adaptation to structural limitations. The very improvisational quality of urban governance, with its negotiated practices and institutional bricolage, makes Indian cities critical sites for reconceptualizing state capacity in the Global South.

By engaging with this messy reality rather than measuring cities against idealized governance models imported from Northern contexts, we uncover a contingent, improvised form of state capacity that paradoxically enables cities to function despite their apparent institutional weaknesses. This perspective helps explain why, despite their fiscal and administrative fragility, complex institutional overlaps, and chronic under-resourcing, local state remain vital arenas of public service delivery, political contestation, and institutional innovation. That makes the paradox all the more intriguing: Indian cities are deeply fragmented yet somehow functional, historically neglected yet increasingly vital. This section therefore asks, provocatively, what ``capacity'' really means in such a context.

\subsection{Historical Neglect and the Late Turn to Urban Reforms}

To fully appreciate the current fragility of the cities—and their paradoxical resilience—one must confront the historical neglect that shaped them. The Indian state effectively abandoned urban governance for the first fifty-odd years after Independence, treating cities as administrative afterthoughts rather than engines of development. Unlike rural development (which saw nationwide programs like community development blocks in the 1950s, Integrated Rural Development in the 1970s, etc.), urban development remained a residual category in national planning. The dominant policy narrative cast urbanization as a byproduct of development, even a necessary evil to be managed, rather than a transformative process to be nurtured.

Two critical junctures mark the belated emergence of urban governance from this residual status. First, the 74th Constitutional Amendment Act of 1992 provided legal recognition to urban local bodies, though without commensurate fiscal or administrative capacity. This half-hearted constitutional acknowledgment created the formal shell of local urban democracy without the substantive powers that might give it life. The second, perhaps more consequential juncture came over a decade later when India launched its first national urban renewal mission. The Jawaharlal Nehru National Urban Renewal Mission (JNNURM), unveiled in December 2005, marked a dramatic shift—the single largest urban development initiative in India's history up to that point. JNNURM explicitly linked central government funding to reform conditions, requiring cities to implement property tax reforms, e-governance, master plans, and other changes to receive infrastructure grants.

Yet this late awakening came after decades of systemic neglect. By 2004, India was already approximately 28\% urban with hundreds of millions living in towns and cities; yet the institutional scaffolding for managing this urban population remained skeletal at best. Even today, combined public spending on urban development by states and Centre hovers around a mere 3\% of GDP, grotesquely disproportionate to what India's urban population ratio would suggest is needed. By contrast, China's decades of high urban investment fueled much smoother service provision, highlighting how this historical neglect has materially shaped India's current urban crisis.

The period since 2004 has witnessed a parade of federal schemes (AMRUT for water supply, Smart Cities Mission, Swachh Bharat for sanitation, and others), each attempting to improve aspects of urban governance and service delivery. These have certainly infused resources and some innovations, yet the scheme-driven approach underscores how the urban local state remains not organically empowered but rather pulled up by external incentives—a dependent rather than autonomous actor in its own development.

This historical trajectory explains why Indian urban governance is as fragmented and frail as it is— because robust local institutions were never a priority until crises emerged. The local state we see today is a product of path dependency: decades of treating cities as administrative wards of states rather than self-governing entities. This historical lens reveals that any evaluation of performance must be tempered by the recognition of this late start—cities have been expected to do more with much less preparation time than other sectors. Yet importantly, it also reveals that improvements, however modest, have occurred despite this historical handicap—suggesting a form of resilience and adaptive capacity that traditional state capacity frameworks struggle to explain.

\subsection{The Local State in Action: Legal Ambiguity and Bureaucratic Agency}

One reason the local state in India functions in a negotiated way is the nature of laws and rules governing it. Urban governance operates in a zone where formal laws are often vague, outdated, or internally contradictory, necessitating interpretation and norm-setting by administrators. The Madhya Pradesh Municipalities Act provides a revealing window into this reality: it creates a formal obligation (provide water) but simultaneously negates any enforceable accountability (no citizen can sue for non-supply). The upshot is that the real obligation is determined by norms, not by the text of law.

This legal ambiguity, rather than merely creating dysfunction, paradoxically creates space for mezo-level bureaucrats to exercise critical discretion. In Bhopal, for example, municipal officials have leveraged this legal ambiguity to extend water delivery to nearly all informal settlements despite no explicit mandate to serve unauthorized areas. The municipal official's interpretation that ``adequate provision'' means doing the best possible with available resources becomes more important than strict legal interpretations. This is governance by bureaucratic agency, not by legislative clarity.

In effect, local governance in India is often government by precedent and practice. Officials learn ``how things are usually done'' from their seniors, and those practices can matter more than what the Acts and rules say. There is considerable room for interpretive maneuver. For instance, if a municipal act says the city ``may'' establish ward committees for citizen participation, whether this happens or not depends on the state's and city's political will. In most states these provisions lay dormant (hence the lack of functional ward committees in many cities), creating a norm of low citizen consultation, despite the 74th Amendment's intent.

Conversely, where reformist bureaucrats have interpreted laws more aggressively, they have occasionally expanded the scope of city action. In Pune, for instance, municipal officials pioneered participatory budgeting practices by creatively interpreting legal provisions on citizen consultation, establishing a model that subsequently spread to other cities. Similarly, Indore's transformation into one of India's cleanest cities occurred not through any revolutionary legal mandate but through bureaucrats leveraging existing provisions while building interdepartmental coordination around solid waste management. These examples illustrate how bureaucratic norms and practice can reshape governance outcomes despite legal ambiguities.

The law thus sends mixed signals to the bureaucracy: ``Do your best, but there's no penalty if you don't succeed.'' In such an environment, what guides action is not legal compulsion but administrative and political culture. If the city or state leadership prioritizes universal water coverage, they will drive the bureaucracy to use its discretionary power to expand services. If not, the clause remains just lofty language. This understanding is a kind of unwritten policy— a tacit agreement that aligns with neither a full laissez-faire approach (they do attempt to provide services) nor a rights-based approach (they stop short of aiming for universality). It is, fundamentally, governance by negotiated settlements: settlements in both the literal sense (informal settlements get water through negotiated means) and the figurative sense (different arms of the state and community ``settling'' for intermediate solutions).

\subsection{Engaging the Local State: From Dysfunction to Transformation}

If fragmented and underpowered city agencies are part of the problem, could they also be part of the solution? It might seem counterintuitive—many policy-makers have bypassed municipal bodies precisely because they appear ineffective. Yet there is growing recognition that sustainable improvement in urban governance requires working with and strengthening local institutions rather than perpetually circumventing them. In other words, the path to transforming city governance runs through the messy local state, not around it.

Odisha's Jaga Mission offers a compelling case. This state-led initiative set out to grant land tenure rights to slum dwellers across cities and upgrade those settlements with basic infrastructure. What makes Jaga Mission particularly remarkable is that it emerged in Odisha—a state traditionally categorized as ``low capacity'' by most development metrics—yet achieved substantial success. Although conceived at the state level, it was implemented largely through existing local governance structures. Municipal officers, ward councilors, and community volunteers in each city were mobilized to survey slum households, identify beneficiaries, and oversee improvements. The state provided the policy framework, funds, and oversight, but on-ground execution relied on city-level actors. The results were remarkable: according to official government sources and independent reports over 2,00,000 families have received secure land titles, and hundreds of slum communities saw tangible upgrades in roads, lighting, sanitation, and more. By empowering slum residents and involving municipal institutions directly in delivering tenure and services, the mission not only improved living conditions but also built confidence and capacity within urban local bodies of a supposedly ``low capacity'' state.

Studies of Basic Services for Urban Poor (BSUP) projects under JNNURM in Chennai and Ahmedabad reveal similar complexities. Coelho, Mahadevia, and Williams (2022) note how ``Chennai's state-led BSUP projects represent a concerted effort to deliver structured housing to the urban poor,'' while Ahmedabad's initiatives ``offer a more flexible model, allowing for incremental upgrades that align with the needs of residents.'' Despite the ``persistent challenge of 'dis-location,' where the relocation of low-income households to peripheral areas disrupts livelihoods and social networks,'' these initiatives demonstrated how urban missions could catalyze meaningful improvements when implemented through local institutions.

The dominant strategy in urban development has become clear: bypass the local state entirely. National programs create parallel structures—special purpose vehicles, state corporations, mission directorates— that sidestep municipal bodies altogether. While this delivers quicker results initially, it reinforces the very weaknesses it seeks to avoid.

The national urban missions reveal this tension perfectly. JNNURM and Smart Cities Mission pumped resources into cities, but local state capacity and priorities determined actual outcomes. Delhi built flyovers for its middle-class constituencies. Mumbai expanded water infrastructure across all neighborhoods, including slums. Bhopal strengthened administrative capacity through management information systems. Lucknow failed to spend funds effectively, concentrating limited investments in elite areas while ignoring wider needs. Despite uniform guidelines and funding, these cities produced radically different results—confirming that local state capacity remains the critical variable.

These missions brought new dynamics to urban governance. International and domestic consulting firms now provide ``private expertise'' at multiple levels, creating what scholars term ``Consultancy Urbanism.'' Dashboard-driven implementation prioritizes metrics over substantive change. Smart Cities abandoned JNNURM's reform agenda entirely, favoring discrete projects instead. Yet these programs created openings for reform that deserve attention. Their limitations—excessive reliance on parastatal agencies, consultant-driven approaches, metric-focused implementation—cannot obscure their significance for understanding how change emerges within constrained institutions.

By not studying these experiences, we miss crucial insights into urban transformation. The local state, though bypassed, remains the essential mediator between national ambitions and urban realities.

\subsection{Rethinking State Capacity: Insights from Indian Cities}

Examining Indian cities as a category forces a fundamental rethinking of dominant models of state capacity, governance, and development. These urban cases defy any simple classification of the state as either capable or failing—instead, they reveal a state that is heterogeneous, internally segmented, and often internally contradictory. This has several profound implications for theory:

\begin{enumerate}
\item \textbf{State Capacity as a Local Variable:} Traditional discussions of state capacity (fiscal capacity, administrative capacity, etc.) often take a country-wide view, or at best distinguish between regions. Indian cities show that capacity can vary not just regionally but organizationally— different arms of the state within the same locality have very different levels of effectiveness. One neighborhood might have reliable electricity (perhaps because a well-run state electricity board manages it), while trash goes uncollected (because the municipal solid waste department is under-resourced and nobody coordinates a solution). Thus, rather than speaking of ``the Indian state'' as having a certain capacity, we see a mosaic of capacities. This underlines the need for more fine-grained theories that look at sub-state units and inter-organizational relations.

\item \textbf{Beyond Principal-Agent: Toward Network Governance Models:} The messy reality of politician--bureaucrat interactions in cities suggests that simple principal-agent models are inadequate. Instead of a linear chain of accountability, Indian urban governance is better seen as a web of relationships—a network governance or multilevel governance system where local and higher-level actors continually negotiate power. This resonates with the idea of polycentric governance (often discussed in the context of metropolitan regions), where overlapping authorities must coordinate (or compete). Indian cities teach us that we should model governance not just as a hierarchy but as a negotiated order. The concept of embedded autonomy itself might be reinterpreted here: perhaps what's needed is an ``embeddedness'' not only in social groups (as Evans meant) but embeddedness across state institutions—for example, city officials embedded in local communities and connected to higher administrative echelons, serving as bridges.

\item \textbf{Democracy and Bureaucracy Reconsidered:} Indian cities challenge the notion that democratic decentralization automatically improves governance. Decentralization theory would suggest that bringing government closer to the people (through elected city bodies) should enhance accountability and service delivery. The Indian case shows that partial decentralization— without real devolution of power and resources—can lead to disillusionment and confusion. It is a cautionary tale that institutions matter more than labels: having a municipal council is not enough if that council cannot command the bureaucracy or finances needed to act. Conversely, the Indian urban experience also suggests that bureaucratic centralization has limits—a top-down approach that marginalizes local knowledge tends to misfire. We see this in the mixed results of one-size-fits-all national urban schemes that ``assumed a level of uniformity that doesn't exist'' across cities.

\item \textbf{The Hybrid State – Formal and Informal Synergy:} Perhaps one of the most original insights from Indian cities is how the formal and informal elements of governance intertwine to produce outcomes. Neither the Weberian bureaucracy model nor pure patron-client model alone fully captures the process; what we see is a hybrid state at work. For example, a formal government program might rely on community volunteers or political intermediaries to actually reach slum populations, essentially blending bureaucratic provision with informal social networks. This is not wholly accounted for in classic models of bureaucratic rationality or clientelism, which often treat informal practices as a deviation. In Indian cities, informality is part and parcel of governance. Recognizing this pushes theorists to move beyond binary thinking (legal/illegal, state/society) toward frameworks that accept a continuum.
\end{enumerate}

These insights extend beyond the Indian context to inform comparative theory-building on state capacity in developing democracies. By moving beyond national-level analyses to examine subnational variation through the lens of conditional autonomy, this framework offers analytical leverage for understanding governance outcomes across diverse political systems characterized by institutional hybridity and contested authority. Similar patterns of fragmented authority, fiscal constraints, and bureaucratic adaptation likely characterize urban governance across much of the Global South, suggesting broader applicability of the conditional autonomy framework developed here.

In making these arguments, Indian cities are not presented as merely failing cases or aberrations. Rather, they are instructive cases that prompt theory-building. They compel us to ask: What does ``state capacity'' really mean when capacity is distributed among different agencies? How do accountability mechanisms function when authority is diffuse? How can equitable public goods be provided in a system that leans on inequitable patronage? The answers gleaned from Indian urban contexts can help reconceptualize state capacity in the Global South as something less about a singular strong vs weak state, and more about the configuration of institutions and relationships that produce developmental outcomes. In other words, Indian cities help shift the discourse from ``Is the state capable?'' to ``Which state institutions are capable, in concert or conflict with which others, and under what conditions?'' – a far richer question.

\section{Conclusion: Opening the Black Box of Urban State Capacity}

Indian cities present a seeming paradox: on the one hand, we see institutional weakness—fragmented authority, low staffing, poor finances, ambiguous laws—and on the other, we find that cities still somehow deliver incremental improvements and maintain a livable order for millions. This chapter has argued that the resolution of this paradox lies in understanding state capacity as an improvised, negotiated practice. The local state in India is not a well-oiled machine, but neither is it a broken engine beyond repair. It is more akin to a jerry-rigged device, patched together and working in unorthodox ways.

Most scholarly analyses of the urban Global South have three tendencies that fundamentally limit our understanding of state capacity: they either focus exclusively on high-level politics (Aijaz, 2008; Kohli, 1987; Kundu, 2003), concentrate on grassroots mobilization and people's movements (Baviskar, 2003; Benjamin, 2000; Shaw, 2007), or reduce everything to neoliberal capital flows (Benjamin, 2007; Ghertner, 2011; Kundu, 2011). As Benjamin (2007) and others have argued, ``new centers of economic, political, and discursive power are emerging in Indian cities as powerful economic and social groups maneuver to achieve their objectives''—yet this neoliberal analytical frame too often obscures the agency and improvisation happening within state institutions themselves.

These dominant approaches neglect the crucial role of the local state bureaucracy where change actually happens—the officials who lay the pipes, build the sewers, and establish waste collection networks. By failing to open this ``black box'' of state functioning, we miss how improvised capacity gets built and incremental change occurs even in constrained circumstances. This dissertation fills this gap by developing a theoretical framework centered on conditional autonomy and meso-level bureaucrats, contributing to comparative politics in three significant ways...he three classic agents of urban change— bureaucracy, politics, and civil society—are all present in Indian cities, yet we have systematically neglected to study how they interact within and through local state institutions.

This perspective carries an important implication: engaging with the urban state is crucial precisely because it is imperfect. If we seek to improve service delivery, reduce inequality, or make cities more sustainable, we cannot bypass city governments and hope for lasting change. Projects that circumvent local institutions (for example, parallel private utilities or ad-hoc task forces) may provide temporary relief but are hard to sustain or scale. In contrast, working through the local state—frustrating though it may be—can lead to institutional learning and transformation. When reformers confront the bottlenecks and push for changes from within (simplifying rules, digitizing systems), they convert bottlenecks into new capacity.

In a broader sense, examining the local state in Indian cities sheds light on state capacity in the Global South. It challenges the binary view of states as either capable or failed. Indian cities show a state that is weak-but-working—a bundle of contradictions where dysfunction and functionality coexist. Infrastructure may be crumbling in parts, yet municipal elections are fiercely contested, indicating faith in the idea of city governance. Bureaucrats may be harried and under-resourced, yet they act as pivotal change-makers when windows of opportunity open. The urban state is both vilified for its incompetence and relied upon for every basic need—truly both a bottleneck and an enabler.

Thus, the dissertation's central claim is intentionally provocative: if meaningful urban transformation is happening at all in India, it is because of the local state, not in spite of it. This runs against a grain of thinking that often credits external actors (private sector, civil society, higher-level governments) for any success and blames the local state for failures. We have seen that external impetus matters (as with national missions), but it ultimately channels through the city institutions. By critically analyzing both the flaws and the potential of urban local bodies, we gain a more nuanced understanding of state capacity. Indian cities teach us that state capacity is often improvised—made on the fly through coalitions of actors, negotiated rules, and creative use of discretion. That is why cities function, even if unevenly.

This perspective provides a new lens for comparative analysis across the Global South. The conditional autonomy framework developed here offers analytical leverage beyond the Indian context, potentially illuminating governance dynamics in diverse settings where formal institutional weakness coexists with pockets of effectiveness. By focusing on how meso-level bureaucrats navigate constraints rather than simply documenting those constraints, this approach contributes to a more agency-centered understanding of governance that can inform both scholarly analysis and practical reforms across developing democracies.






\section{Bhopal and Lucknow: Similar Institutions, Radical Outcomes}

Having established how conditional autonomy emerges within India's fragmented urban governance structure, we now turn to examining the empirical case studies that illuminate these theoretical propositions.

\textbf{{Scope and Argument in Brief}}

The argument in this dissertation is that the key to understanding change in low capacity, patronage setting lies in the relationship between politics and bureaucracy.
The governance framework identifies four modes based on political incentives (programmatic vs. patronage) and bureaucratic autonomy (autonomous vs. instrumentalized):
\begin{itemize}
  \item \textbf{Programmatic:} Autonomous bureaucracy + programmatic politics = high performance
  \item \textbf{Collusion:} Instrumentalized bureaucracy + patronage politics = governance failure
  \item \textbf{Conditional Autonomy:} Autonomous bureaucracy + patronage politics = negotiated effectiveness
  \item \textbf{Bureaucratic Minimalism:} Weak bureaucracy + low patronage = minimal governance.
\end{itemize}

Figure 12 is repeated from section 2 to illustrates these four governance contexts and locates the dissertation's case studies: the city of Bhopal exemplifies Conditional Autonomy, whereas Lucknow typifies Collusion.

\begin{figure}[htbp]
  \centering
  \caption{Governance typology matrix}
  \label{fig:governance_matrix}
  % You'll need to add your figure here
\end{figure}

To rinstate the theory the central thesis is that meso-level bureaucratic autonomy under patronage politics (``conditional autonomy'') enables substantially better public service outcomes than the typical patronage-collusive equilibrium. In essence, even in a patronage political context, if mid-tier officials – the ``connective tissue'' of the state – enjoy conditional discretion, they can coordinate complex tasks and push for broader provisioning of services. Conversely, when bureaucrats are tightly instrumentalized by patronage networks, governance deteriorates into piecemeal favors and fragmented services. Crucially, this effect is most pronounced for non-targetable public goods – those that are networked and indivisible (e.g. citywide water systems, sewerage, electricity grids) – which require extensive coordination and cannot be sliced into local patronage gifts. By contrast, targetable (divisible) goods – those deliverable in localized units (e.g. street paving, neighborhood trash pickup, handpumps or public toilets) – can still be provided under patronage (though unevenly). The proof of bureaucratic coordination, therefore, lies in the delivery of complex, networked services: no single politician or frontline worker can singly supply a citywide water network or integrated sanitation system. If we observe such goods being extended broadly, it signals an underlying bureaucratic mechanism at work. In short, to understand divergent outcomes under similar institutions, we must examine how politician– bureaucrat relations configure the delivery of different types of goods.

How does ``conditional autonomy'' actually translate into better services? The causal pathway begins with the meso-level bureaucrats – those in intermediate ranks who bridge high-level policy and street- level implementation. Under conditional autonomy, these officials leverage their professional norms and coordinating authority to pursue public goods that politicians with short-term electoral goals might neglect. They act as institutional intermediaries, aligning technical imperatives with political incentives. Given a degree of freedom, bureaucrats can channel resources to citywide projects, synchronize efforts across departments, and insist on rule-bound delivery even while accommodating some political input. In Bhopal, for example, a cadre of state-appointed municipal officers strategically balanced the politicians' patronage demands (ensuring, say, that each councilor got some local projects) with an overarching push for citywide improvements in water and sanitation. This contrasts with Lucknow, where bureaucrats, lacking autonomy, became purely reactive instruments of their political masters – unable to plan beyond parochial interests or electoral cycles.

The nature of the public good conditions this dynamic. For targetable/divisible goods (e.g. routine trash collection, minor road repairs), even a patronage-driven administration can achieve moderate coverage – albeit selectively. Patronage politicians prefer such goods because they can be distributed to specific loyal communities as rewards. Indeed, under collusion we expect a pattern of localized patronage: services concentrated in vote-bank areas and scant elsewhere. A bureaucratically autonomous regime, on the other hand, tends toward universalized delivery of even targetable goods – extending access broadly as a matter of policy, not favor. The real test of governance comes with networked, non- targetable goods.

These require infrastructural complexity – pipes, pumps, grids, treatment plants – that must work as an integrated whole. Such goods cannot be meaningfully confined to one neighborhood or traded in small units. As a result, a patronage-collusive system struggles with them: politicians see little electoral payoff in investments that are citywide, long-term, and cannot be allocated exclusively to supporters Bureaucrats in a collusive system also lack the mandate to coordinate the many components and agencies involved. The outcome is complete fragmentation – patchwork provision, half-built networks, chronic under-service. By contrast, under conditional autonomy, bureaucrats can incrementally push these large-scale projects forward – budgeting for network expansion, enforcing technical standards, and knitting together disparate efforts into a cohesive plan The outcome I term ``incremental universalism'': a significant expansion of networked services (beyond what patronage alone would deliver), though perhaps not perfect or fully institutionalized. Empirically, this framework predicts that differences will be modest for simple, targetable goods but dramatic for complex, non-targetable goods This is exactly what we observe in the comparative cases of Bhopal and Lucknow.

Consider service delivery in Bhopal and Lucknow to illustrate these dynamics. Trash collection, a targetable good, enables selective routing—patronage officials can provide daily pickup to loyal neighborhoods while neglecting others. Water supply, a networked good, requires coordinated investment in infrastructure spanning all areas.

Data from a 2022 (Table 4) citizen survey confirms theoretical expectations: Lucknow's collusive system delivers for trash—49\% of households receive daily collection versus Bhopal's 76\%. More starkly, 12\% of Lucknow residents report no service (concentrated in lower-caste areas) compared to virtually none in Bhopal.

For water—the networked good—disparities intensify: Bhopal achieved 87\% household connections versus Lucknow's 72\%, despite similar baselines twenty years prior. Most tellingly, Bhopal delivered 95\% water access to marginalized slum areas, while 25-30\% of Lucknow's population remains without formal connections, creating a stark patchwork of service provision.

Analysis of additional goods reinforces this pattern. For Non-targetable goods: For electricity service, Lucknow residents face an 82.3 percentage point lower probability of reliable power compared to Bhopal. Sanitation access shows a 13.3 percentage point reduction in Lucknow.

For Targetable goods like quality drain (naali), Lucknow experiences 11.7 percentage points higher probability of waterlogging problems. The exception is road quality.

Results confirm the framework: non-targetable services (electricity, drainage) show massive gaps favoring Bhopal, while more targetable services (roads) present mixed evidence, reflecting how conditional autonomy most impacts complex infrastructure requiring coordination.

\begin{table}[htbp]
\centering
\caption{Public service outcomes in Bhopal vs. Lucknow. The conditional autonomy context of Bhopal achieved broadly inclusive services high coverage of trash and water, low exclusion, whereas Lucknow's patronage-driven context produced fragmented outcomes. Networked goods (water, electricity) show especially large gaps, reflecting the coordination challenge of non-targetable services.}
\begin{tabular}{lcc}
\hline
\textbf{Public Service (circa 2021–22)} & \textbf{Bhopal (Conditional Autonomy)} & \textbf{Lucknow (Patronage Collusion)} \\
\hline
Households with daily trash collection & 76\% (citywide, inclusive) & 49\% (localized pockets) \\
\hline
Households with no trash service & 0\% (nearly universal) & 12\% (many slums unserved) \\
\hline
Households with piped water on premises & 87\% & 72\% \\
\hline
Households with reliable electricity (no power cut) & Vast majority ($\approx$80–90\%) & Tiny minority ($<$10\%) \\
\hline
\end{tabular}
\label{tab:services}
\end{table}

To reiterate the key theoretical insight: service delivery outcomes generate self-reinforcing feedback loops that reshape political-institutional equilibrium. In Bhopal, near-universal service provision (76\% trash collection, 87\% water access) established new norms of universality, making service retraction politically costly and institutionalizing bureaucratic gains. Conversely, Lucknow's fragmented delivery (49\% trash, 72\% water) normalized low expectations and patron-client dependencies. This created divergent trajectories—Bhopal's virtuous circle of incremental reform versus Lucknow's vicious circle of institutional stagnation—demonstrating how conditional autonomy alters governance through iterative service improvements.

\section{Research Design and Methods}

To investigate these propositions, the study employs a comparative subnational research design with both cross-city and within-city analysis. The two primary cases – Bhopal and Lucknow – were selected as a controlled comparison using the Most Similar Systems Design.

Both are second-tier metropolises and state capitals located in India's northern Hindi heartland (often dubbed the ``BIMARU'' belt for its historically weak development).

Table 5 outlines a few key similarities in their profiles. Despite Bhopal being in Madhya Pradesh and Lucknow in Uttar Pradesh, the two cities share common initial conditions: both were mid-sized cities ($\sim$1.8–2.8 million population) with similar literacy rates and socio-economic compositions; both had a legacy of indirect rule (former princely states integrated after independence) and adopted identical municipal institutions under India's 74th Constitutional Amendment (e.g. directly elected mayors, city councils, and state-supervised municipal commissioners) . Through the 2000s, both cities came under the same national urban programs and were ruled by the same national party (the BJP) at both state and municipal levels. In short, formal structures and macro-political contexts were held constant, making the divergence in outcomes a puzzling result of meso-level factors rather than structural disparities. Bhopal and Lucknow thus form a ``most-similar'' pair in which one can trace how different bureaucratic responses under similar pressures led to very different governance outcomes.

\begin{table}[htbp]
\centering
\caption{Profile of Bhopal and Lucknow. The cities share demographic and institutional baselines but differ in administrative arrangements. Notably, Bhopal's governance is more integrated and bureaucratically professionalized, whereas Lucknow's is more fragmented and politicized. These differences became salient after 2010, when the cities embarked on divergent reform paths.}
\begin{tabular}{lcc}
\hline
\textbf{City} & \textbf{Bhopal (Madhya Pradesh)} & \textbf{Lucknow (Uttar Pradesh)} \\
\hline
Population (2011) & $\sim$1.80 million & $\sim$2.82 million \\
\hline
\% Literate & $\sim$83\% & $\sim$86\% \\
\hline
\% Slum Population & $\sim$27\% (PCA Census data) & $\sim$13\% (PCA Census data) \\
\hline
Governance Structure & Unified Municipal Corporation (water, waste, etc. under city) & Fragmented agencies (e.g. state-run Water Board separate from city) \\
\hline
Municipal Cadre & Yes – Dedicated urban service cadre; stable core staff & No – Posts filled by general cadres; frequent turnover \\
\hline
Commissioner Tenure & 24 months avg.; always IAS officers (elite civil service) & 19 months; often non-IAS until recently (mezo level) \\
\hline
Political Control & BJP-dominated (state \& city) since 2003 & BJP-dominated (state \& city) since 2017 \\
 & & BJP dominated city since 2006 \\
\hline
\end{tabular}
\label{tab:profile}
\end{table}

Beyond the city selection, the research design incorporates a careful within-city sampling and process-tracing strategy. A large-N household survey (N$\approx$3,754) was conducted in 2022 across both cities to systematically measure service delivery outcomes and citizen experiences. Ward selection for the survey followed a stratified random approach: in each city, wards (the lowest urban administrative units) were chosen to represent inner-city and peripheral zones, varying built-up environments, and a mix of communities including Scheduled Castes, Scheduled Tribes, and Muslims. This ensured that the sample captures the full social and spatial diversity of each city – from old dense quarters to new informal settlements. The household survey gathered data on access to services (water, trash, toilets, electricity, roads), frequency and reliability of provision, and citizens' modes of obtaining services formal requests, political networks, etc. It also measured political participation, allowing analysis of whether citizens' engagement yields different returns under different city regimes. Survey results quantitatively confirm the theorized patterns (as shown earlier in Table 1): Bhopal decisively outperforms Lucknow on most indicators, especially non-targetable goods, and political networking is far less necessary to get basic services in Bhopal than in Lucknow.

The research employed a controlled within-city comparative design. Four wards were systematically selected in each city based on a paired sampling methodology (Table 6,7). This approach created matched neighborhood pairs across Bhopal and Lucknow - for instance, comparably low-income slum-dominated wards in both cities. Such controlled-pairing allows ``like with like'' comparison, isolating governance effects from demographic variation. This micro-level matching strengthens causal inference by attributing observed service delivery differences to governance models rather than local demographic characteristics.

\begin{table}[htbp]
\centering
\caption{Field Sites in Lucknow}
\begin{tabular}{lll}
\hline
\textbf{Ward} & \textbf{Trash Collection} & \textbf{Rationale} \\
\hline
Aliganj & ``Contractor'' and city empaneled private concessionaire & Older Development - Salaried Class, Middle Class, Planned Colony \\
\hline
Gomti Nagar & City empaneled private Concessionaire & New Development - Upper Middle Class, Commercial establishments \\
\hline
Hazratganj & ``Contractor'' & Old Development - Significant Muslim population \\
\hline
Transport Nagar & ``Contractor'' and city empaneled private concessionaire & New Development - Migrant population, Dalits and Muslims from Bihar and Eastern Uttar Pradesh \\
\hline
\end{tabular}
\label{tab:lucknow}
\end{table}

\begin{figure}[htbp]
\centering
\caption{Map of Bhopal City by Wards – Field Sites Highlighted in Color}
% Insert your figure here
\label{fig:bhopal_map}
\end{figure}

\begin{table}[htbp]
\centering
\caption{Field Site in Bhopal}
\begin{tabular}{lll}
\hline
\textbf{Ward} & \textbf{Trash Collection} & \textbf{Rationale} \\
\hline
Bhanpur & Legacy waste site & Older Development - had a 30-year-old legacy site. Comparable to Shivari site in Lucknow. \\
\hline
Maharana Pratap Nagar & Bhopal City Government & New Development - Middle and upper middle Class, Commercial establishments. Comparable to Gomti Nagar in Lucknow. \\
\hline
Bawadiya kalan & Bhopal City Government & New Development - Migrant population, Tribals and Dalits from MP, Rajasthan and Chhattisgarh. Comparable to Transport Nagar in Lucknow. \\
\hline
Teela Jamalpur & Bhopal City Government & Old Development - Significant Muslim population \\
\hline
Indira Gandhi & Bhopal City Government & Informal colony - slum. To test Bhopal City Government's capacity to cover informal settlements. \\
\hline
\end{tabular}
\label{tab:bhopal}
\end{table}

The study combines statistical analysis with qualitative process-tracing based on fieldwork conducted from 2021-2023. Data collection included:

\begin{itemize}
\item In-depth interviews with bureaucrats (elite, meso-level, frontline), politicians (city and state), and civil society actors in both cities
\item Archival research on policy decisions and reforms since the 1990s The analysis focuses on critical junctures:
\item 74th Constitutional Amendment implementation (municipal devolution)
\item JNNURM rollout in the 2000s (infrastructure funds tied to reforms)
\item City-specific initiatives (e.g., Bhopal's municipal services cadre creation, 2011)
\end{itemize}

Process-tracing reveals how these cities diverged at key moments. Bhopal embraced JNNURM reforms to modernize services; Lucknow's agencies resisted change. Bhopal officials secured political buy-in by framing technical programs as political wins, gaining operational autonomy. Lucknow's bureaucrats faced tight political control. This comparative process-tracing establishes the causal chain: initial conditions → governance mode → implementation strategies → outcomes.

\begin{figure}[htbp]
\centering
\caption{Research design flowchart. The study proceeds through case selection and pairing, data collection, and comparative analysis. Bhopal and Lucknow were chosen as most-similar cases (both state-capital cities with comparable profiles) differing in hypothesized governance mode. Within each city, wards and neighborhoods were sampled to ensure diverse and comparable sub-units. Mixed-methods evidence – including a 2022 household survey and 2021–23 interviews/observations – was then analyzed to trace how each city's bureaucrat–politician interactions affected service delivery outcomes.}
% Insert your figure here
\label{fig:flowchart}
\end{figure}

This multi-method design combines breadth and depth. The household survey (n=3,754) provides generalizable evidence of service delivery differences and correlations with political/bureaucratic factors. Qualitative process-tracing explains why these differences emerged through interviews and archival research.

Triangulating methods strengthens validity. The analysis shows state capacity emerges from meso-level interactions, not uniform national/state characteristics. Despite identical formal institutions, Bhopal evolved toward inclusive service delivery while Lucknow remained clientelistic.

The findings reveal a ``conditional autonomy'' zone between Weberian ideals and patronage collapse, where partial reforms yield substantial public goods improvements. This demonstrates how pockets of effectiveness form within adverse environments and how empowering meso-level bureaucrats transforms nominal institutions into functional governance. The comparative analysis suggests incremental institutional change can build effective state capacity even in patronage-heavy systems.

\section{Patronage vs. Incremental Reform – Systematic Comparison and Causal Mechanisms}

The stark divergence between Bhopal and Lucknow can now be distilled into a comparative analysis of their underlying governance models. In brief, Bhopal followed a programmatic path of incremental reform, whereas Lucknow remained stuck in a collusive patronage trap (Both cities had comparable formal mandates and faced similar resource constraints, yet their trajectories could not be more different. This section examines why Bhopal succeeded in gradually strengthening institutions and service delivery, while Lucknow entrenched a clientelistic status quo. The contrast illuminates key causal mechanisms – in bureaucratic behavior, political incentives, and center–local relations – that can make or break urban governance in low-capacity settings.

\begin{figure}[htbp]
\centering
\caption{Causal Model – Conditional Autonomy and Public Goods}
% Insert your figure here
\label{fig:causal_model}
\end{figure}

Figure 13 is the replication of my theoretical model for the cases. This model maps how weak local states navigate governance reforms through critical decision points. The starting point is external or internal reform pressure—in India, the 74th Constitutional Amendment Act (CAA) served as the first major push for urban decentralization, followed by JNNURM (Jawaharlal Nehru National Urban Renewal Mission) which tied infrastructure funding to governance reforms. These pushes create decision points where political elites must choose whether to grant bureaucratic autonomy.

\begin{itemize}
\item \textbf{The Decision Point:} When reform pressure arrives, political elites face a choice: support bureaucratic autonomy or maintain tight control. This decision determines which path a city takes. In the case of Lucknow elites chose to not extend autonomy and in Bhopal they did.
\item \textbf{Left Path (Lucknow):} Without political support, bureaucrats remain controlled by politicians and power brokers. This leads to collusive behavior where bureaucrats align with political interests for survival. Service delivery becomes selective—targetable goods flow to patronage networks while universal services are neglected. The outcome: localized patronage perpetuates fragmentation.
\item \textbf{Right Path (Bhopal):} With political support, mid-level bureaucrats gain conditional autonomy—discretionary space to coordinate across departments and innovate. This triggers a virtuous cycle: bureaucrats form teams, introduce new systems, learn from projects, and build institutional memory. Service delivery improves—targetable goods become more universal, non-targetable services expand incrementally. The outcome: institutional embedding leads to universalism.
\item \textbf{The Bidirectional Risk Mechanism:} This is the model's key dynamic. Once autonomy is granted, politicians could revoke it, but improved services create citizen expectations. Politicians then face electoral backlash if they reverse reforms—this risk calculation stabilizes autonomy. The relationship works both ways: bureaucrats know autonomy is conditional, motivating performance; politicians know reversal carries costs, encouraging persistence.
\item \textbf{First Step:} In case of conditional autonomy conditions, incremental progress is made. If the support sustains capacity builds in the local state.
\item \textbf{The Feedback Loop:} As service delivery improves, cities attract external funding and support from higher government levels and institutions. This reinforces bureaucratic capacity and autonomy, creating a self-sustaining reform trajectory.
\end{itemize}

Each new reform push tests this system—cities may shift paths if political calculations change. The model reveals how weak states can transform governance through strategic autonomy grants at critical junctures. National representative data supports the model and its claim. Figure 14 below maps change in public service over time for Bhopal and Lucknow with India (urban). We observe modest increase in the decades of 2000 and then once these reforms institutionalize, there is a spike in delivery of public goods.

\begin{figure}[htbp]
\centering
\caption{Authors calculation using NSS data for Bhopal Lucknow and India}
% Insert your figure here
\label{fig:nss_data}
\end{figure}

\subsection{Bhopal's Incremental Capacity-Building}

The process-tracing revealed that a major factor in Bhopal's success was the presence of bridging actors in the bureaucracy who could intermediate between political pressures and technical needs. These state-appointed municipal officers in Bhopal effectively became ``dual agents'' – accountable upward to state government but deeply embedded in local administration. They cultivated what Peter Evans might call embedded autonomy, maintaining enough independence to pursue coherent policies while staying connected to on-the-ground realities. This was facilitated by MP's institutional choice to have a dedicated municipal service cadre and to embed reform projects within the city government rather than creating parallel structures. For example, during the JnNURM water supply upgrade, the project unit was housed in Bhopal's municipal offices and led by Madhya Pradesh Municipal Service officers. This meant that those implementing the project were not outsiders but part of the local team, familiar with the city's quirks and invested in its success. It also meant they had to navigate local politics daily. Instead of seeing political interests as an obstacle, Bhopal's bureaucrats learned to accommodate and shape them. They actively built bridges – informing councilors and MLAs of project benefits, adjusting technical plans to address local complaints, and finding compromises so that political actors could claim some credit. As noted, Bhopal's officials found ways to include local contractors often tied to politicians in procurement, but insisted on transparent bidding for at least a portion of works. They created informal norms of ``no outright losers'' – stakeholders who might have blocked a reform were instead co-opted, but on terms that still advanced the project. Over time this generated a reservoir of trust. Politicians saw that bureaucrats were not there to undercut them; bureaucrats saw that politicians, once convinced of a project's value, would let them implement it with minimal interference. This tacit bargain – you help us deliver development, we won't completely sideline your patronage – underpinned Bhopal's iterative progress. Each successful project built confidence and a precedent for the next. By institutionalizing joint problem-solving in this way, Bhopal avoided the fate of many cities where reforms start with fanfare and then fizzle due to turf wars. Instead, Bhopal created an ``adaptive governance'' culture: reformers incrementally expanded services by working with political incentives, not against them, and by cleverly packaging technocratic improvements as also being in the politicians' interest. This culture took years to solidify. But once established, it meant the city could marshal a united front when opportunities arose (like new funding or programs). The modest but steady gains in Bhopal attest to the power of such alignment.

Several causal mechanisms specific to Bhopal stand out. First is the presence of a reform coalition early on. Key individuals – a supportive Chief Minister in the 1990s, a capable mayor or municipal commissioner at critical times – acted as champions for the city's autonomy. They forged a mutual dependence: local politicians needed bureaucrats to fulfill development promises, and bureaucrats needed political cover to make administrative changes. This reciprocity discouraged open conflict and encouraged ``getting to yes'' on reforms. Second, positive feedback loops took hold. Minor successes (e.g. improving one ward's water supply) were publicized and generated citizen goodwill, which in turn gave political leaders an electoral incentive to allow further technocratic work. Community satisfaction surveys began reflecting these improvements, shifting the narrative from ``Bhopal is hopeless'' to ``Bhopal is improving,'' which bolstered reformers' morale. Third, Bhopal benefited from not having an alternative patronage pipeline. Unlike Lucknow, which as a politically central city could expect bailouts or special favors, Bhopal was relatively peripheral (MP is not a powerhouse in national politics). Thus, when union funds or donor projects became available, Bhopal's elites knew ``perform and get the funds, or get nothing.'' This created a shared urgency and disciplined the local decision-makers to meet program conditions. In Lucknow, by contrast, there was always the hope of obtaining resources through lobbying or partisan channels, reducing the pressure to reform. Fourth, institutional design mattered: MP's decision to grant cities more fiscal and administrative power (and to nurture a professional urban cadre) paid dividends. Bhopal's mayors wielded meaningful authority and took pride in their city, rather than serving merely ceremonial roles. A consultant observed that ``cities in Madhya Pradesh were not pushovers – the Mayors knew they were powerful and maintained independence from state leadership.'' This local assertiveness kept state-level meddling in check to an extent, allowing bureaucrats to focus on local priorities.

\subsection{Lucknow's Collusive Patronage Trap}

In stark contrast stands Lucknow, where the same mechanisms turned corrosive. Lucknow's governance operated through what one might call a ``networked patronage'' model. Political and bureaucratic elites formed tightly interwoven alliances aimed not at solving collective problems, but at maintaining their own networks of support. The immediate, tangible benefits of this system (jobs for supporters, lucrative contracts for favored businesses, informal payments flowing upward) created strong incentives to preserve the status quo. Meanwhile, the costs – deteriorating infrastructure, inequitable services – were diffused among the general public and accrued slowly, often hidden until crises hit. This misalignment of incentives meant that there was no internal push for reform: those with power were profiting from dysfunction, and those suffering had little voice. Lucknow's bureaucrats, far from acting as intermediaries for the public good, became clients of political patrons. A stark illustration is the massive politicization of municipal hiring, which filled the ranks of the city with employees who owed their jobs to a patron. This selective inclusion did in fact extend some benefits to marginalized communities for example, many Valmiki Dalits got government jobs under Mayawati's tenure. However, it also entrenched a norm that loyalty trumped merit, and that the purpose of public institutions was to distribute rents rather than services. Any bureaucrat trying to prioritize performance or long-term planning in Lucknow faced a daunting barrier: a labyrinth of vested interests with veto power at every turn. As we saw, even seemingly straightforward upgrades (like installing water meters to reduce waste) were torpedoed by powerful ministers to preserve patronage-based billing loopholes. The water sector was especially fragmented – the legally autonomous water utility was bypassed and undermined by state agencies , and local knowledge was ignored. An observer noted that those most familiar with Lucknow's water needs (the utility engineers) ``were excluded from critical decisions'', leading to plans that made little practical sense. It's a textbook case of how informal power (``shadow state'') can paralyze formal institutions. When DFID and ADB came with money, rather than unite to leverage it, Lucknow's actors fragmented further – each trying to capture a piece of the pie (The result: projects collapsed, and no one was held accountable.

The causal logic in Lucknow is a negative mirror of Bhopal's positive feedback loop. In Lucknow, each failed reform further convinced elites that real change was impossible (or not in their interest), thus reinforcing their reliance on patronage fixes. When one private trash contract failed, there was no move to build municipal capacity; instead, officials simply awarded another contract to a different crony firm, hoping for a different outcome (which again failed) (This is akin to a vicious cycle of learned helplessness in governance. Another mechanism at play is the fragmentation of authority: whereas Bhopal cultivated a degree of alignment across levels, Lucknow became a story of each level – city, state, even national – working at cross purposes. The metaphor of a three-legged race with partners running in different directions is apt (City officials in Lucknow often found their initiatives blocked by state-level bosses; state politicians found their grand schemes sabotaged by local cadre incompetence or citizen resistance. The lack of a cohesive ``reform coalition'' meant that even well-intentioned individuals were quickly isolated and neutralized. Many capable IAS officers avoided Lucknow postings or left quickly, finding the environment unconducive to professional accomplishment. Those who remained adapted by becoming part of the patronage machine. In essence, Lucknow's system had a remarkable capacity to resist change. It co-opted new policies turning, say, a public-private partnership into just another patronage outlet or waited out reformist officials until they were transferred. This path dependence yielded a stable equilibrium – stable in political terms, that is, even as service delivery stagnated. Breaking out would have required an extraordinary coalition and incentive realignment that, absent external shocks, never materialized. Thus, Lucknow's case underscores how patronage-driven governance can nullify even well-funded or well-designed reforms, by bending them to serve the existing network or by creating enough administrative confusion that nothing effectively changes.

To systematically compare, we distill a few key differences:

\begin{itemize}
\item \textbf{Bureaucratic Autonomy vs. Subservience:} In Bhopal, mid-level bureaucrats attained conditional autonomy – they weren't free agents, but they had enough stability and backing to act in the public interest when opportunities arose. n Lucknow, bureaucrats were largely subservient to political patrons, with rapid turnover ensuring compliance over initiative. A stable bureaucratic tenure in Bhopal meant institutional memory and steady capacity growth; Lucknow's revolving-door postings meant institutional amnesia, with each new commissioner ``reinventing the wheel'' (or abandoning it).
\item \textbf{Coalition-Building vs. Fragmentation:} Bhopal's governance forged a working coalition among the mayor, municipal staff, state officials, and even donors – aligning incentives (at least partially) toward common development goals (Lucknow's governance was marked by fragmentation and infighting – city politicians, state ministers, local contractors, and different bureaucratic fiefdoms each guarded their turf (Where Bhopal minimized veto points through informal understandings, Lucknow seemed to maximize veto players, making coordination for the collective good nearly impossible.
\item \textbf{Problem-Solving Norms vs. Rent-Seeking Norms:} In Bhopal, a norm evolved that valued pragmatic problem-solving – e.g. if a sewer line was blocked, the question was how to fix it with available means (and maybe get political credit along the way). In Lucknow, the dominant norm was who gets to control the contract or job associated with the sewer, rather than ensuring the sewer actually works. These norms become self-reinforcing. Bhopal's officials took pride in service improvements (several interviewed officers spoke of projects they initiated that continued after their transfer, indicating institutionalization). Lucknow's officials, in contrast, often treated service delivery as someone else's problem – their focus was on maintaining the network that kept them in power.
\end{itemize}

Ultimately, Bhopal created a virtuous cycle, Lucknow a vicious cycle. Bhopal's gradual, negotiated improvements emboldened reformers and expanded what was politically possible next time (Lucknow's fragmented, patronage-driven efforts neutralized each attempted reform, teaching actors that broad changes were futile and reinforcing patron-client relationships instead (These pathways had direct consequences for citizens' lives, as the next section shows. Before turning to those outcomes, it is worth noting that neither city's path was preordained – small decisions at critical junctures accumulated into very different state capacities. Bhopal and Lucknow demonstrate that governance arrangements, and specifically the presence or absence of coordination mechanisms, are the decisive factor shaping urban service trajectories.

\section{Conclusion}

The divergent trajectories of Bhopal and Lucknow underscore how sub-national incentives—rather than formal structures or resource endowments—shape state capacity in low-income settings.
Beginning from similar starting points in India's BIMARU heartland, Bhopal leveraged political-bureaucratic alignment to build incremental ``islands of effectiveness,'' whereas Lucknow became trapped in a patronage equilibrium that kept basic services unreliable. The comparison dispels deterministic views of urban under-development: even weak states can improve when institutions, tenure rules, and political incentives converge on programmatic goals.

Lucknow illustrates a self-reinforcing collusive trap: politicians and bureaucrats exchanged jobs, contracts, and favors for support, crowding out systemic improvements. Because benefits were distributed as individualized patronage, citizens relied on brokers rather than collective voice, bureaucrats were co-opted, and external reform attempts—water metering, garbage privatization— were blocked whenever they threatened the network's rents. Resources flowed, but capacity stagnated.

Bhopal offers the counterfactual. A stable cadre of mid-level municipal officers, backed by reform-minded politicians and external partners, created pockets of competence that gradually expanded city-wide. Coordination among city, state, and donors accelerated infrastructure roll-out; institutional memory accumulated through longer tenures; and universalist service expansion lifted the baseline for all, including the poor. The case affirms that modest, well-sequenced reforms can compound into transformative change without wholesale regime overhaul.

Policy lessons follow. Empower mid-level ``connectors'' with secure tenures and performance incentives; pursue decentralization that grants genuine fiscal and administrative autonomy; and align multi-level actors around shared benchmarks. Anti-corruption drives or privatization schemes will backfire if they ignore underlying incentive structures. Ultimately, effective urban governance is less a technical puzzle than a political one: shifting patronage-based calculus toward performance-based rewards is the fulcrum on which lasting capacity turns.

\section{Dissertation Chapter Organization}

\begin{tabular}{|l|l|l|l|l|}
\hline
\textbf{Number} & \textbf{Title} & \textbf{Method} & \textbf{Question} & \textbf{Contribution to the Argument} \\
\hline
1 & Introduction & & & Sets up the research question and the theoretical framework, providing an overview of the dissertation's structure. \\
\hline
2 & Patronage, Networks and Urban India & Quantitative & Is there a relationship between bureaucrats and nontarget able goods? Do networks matter to access goods? & This section sets up the empirical puzzle: despite comparable formal structures, Bhopal delivers public goods more effectively than Lucknow. In low-capacity Indian cities, citizens rely on informal networks to access services—a pattern especially pronounced in Lucknow. Although these citizen-politician-bureaucrat ties marginally raise individuals' chances of receiving benefits, they leave overall city-level provision weak. Crucially, the analysis shows that networks do not improve non-targetable goods, reinforcing the causal claim that greater conditional bureaucratic autonomy—not patronage networks—drives better service delivery. \\
\hline
3 & Garbage Management & Qualitative: 4 field sites, interviews with IAS and State Services officials & How do bureaucrats manage garbage services in Bhopal compared to Lucknow? & Shows that Bhopal's bureaucrats operate with greater autonomy and internal coherence, enabling effective garbage management. This chapter provides mechanism to understand almost equal delivery of trash delivery in Bhopal and Lucknow extending the empirical results in chapter 2. \\
\hline
4 & Water Management & Qualitative: 4 field sites, interviews with IAS and State Services officials & How do bureaucrats manage water services in Bhopal compared to Lucknow? & It underscores the autonomy of Bhopal's bureaucrats and their capacity to coordinate with external actors, fostering more participatory and effective water management. It also documents how the meso-level bureaucracy has incrementally institutionalized water services over time. \\
\hline
5 & Conclusion & Quantitative: 3,000 sample survey of citizens & Do citizens in Lucknow and Bhopal receive differentiated services, and is bureaucracy the key intervening variable? & Tests the theory's claims by comparing citizen experiences in both cities, confirming that bureaucratic autonomy and coherence are central to effective service delivery. \\
\hline
\end{tabular}



















